\chapter{Biomechanics}
\label{vol06:ch06}

\epigraph{The body is a machine, but a machine of a very peculiar
sort: one that builds itself, repairs itself, controls itself, and
that knows when to stop.}{Y.\ C.\ Fung, \emph{Biomechanics: Mechanical
Properties of Living Tissues} (2nd ed., 1993), introduction
\cite[p.~1]{text:fung-biomechanics}.}

\chapmeta{Half-life: foundational. Archetypes: balance, scaling, failure. Exercise target: 30. Project track: physical experiment plus analysis (instrumented vertical jump with smartphone accelerometer, peak knee force estimation by inverse dynamics). Hazard class: low (the jumping subject should be in normal health; the instrumentation is non-invasive). Reader path default: standard, with sections explicitly tagged. Named cases: none required; orthopaedic-implant and musculoskeletal failures are described as mechanism classes.}

\noindent
Biomechanics applies the mechanics of Volume III to living systems.
The chapter develops bones, joints, and muscles as engineered
structures; locomotion across animals; the cardiovascular system as
a pulsatile hydraulic network; the respiratory system as a
gas-exchange device; the mechanics of single cells; the constitutive
modelling of soft tissues; and the failure modes of musculoskeletal
tissues and orthopaedic implants. The level of detail is engineering
working depth: enough to design a prosthesis or to interpret a
gait-analysis report, not enough to substitute for an orthopaedic
surgery residency. The opening sections treat the body's structural
hardware (bones, joints, muscles, tendons) at the same level Volume
III treats columns, beams, and bearings, and the cardiovascular and
respiratory sections treat the body's principal fluid-mechanical
systems at the same level Volume IV treats pipes, pumps, and heat
exchangers. The closing sections on cell mechanics, soft-tissue
constitutive models, and orthopaedic implants make explicit the
chain from the molecular-mechanics signals an integrin transduces to
the load-bearing performance of a $20$-year hip prosthesis.

\section{Bones, joints, muscles as structures}\pathtag{core}

A bone is a composite material; a joint is a low-friction bearing; a
muscle is an active actuator. The three components together form the
musculoskeletal system, the body's mechanical infrastructure.

\engterm{Bone composition and structure}. Bone consists of an organic
phase (mostly type-I collagen, $\sim 30\,\%$ by mass) and a mineral
phase (mostly hydroxyapatite, $\sim 60\,\%$) plus water ($\sim 10\,\%$).
The collagen provides toughness; the mineral provides stiffness; the
composite has Young's modulus of $\sim 15\,\text{GPa}$ along the
fibre direction and compressive strength of $\sim 200\,\text{MPa}$
in cortical bone. Cortical (dense) bone forms the outer shells of
long bones; trabecular (spongy) bone fills the interior with a
network of struts whose orientation aligns with the principal
stress directions. The trabecular orientation is the canonical
example of Wolff's law (1892): bone adapts its internal structure to
the prevailing mechanical loads
\cite[ch.~1]{paper:v6c06-wolff-1892}. The adaptation is mediated by
osteocytes embedded in the bone matrix that detect strain and signal
to osteoblasts (bone-building cells) and osteoclasts (bone-resorbing
cells); Frost's mechanostat formalises the feedback as a strain
setpoint at the osteocytes
\cite[p.~74]{paper:v6c06-frost-1987}.

Figure~\ref{fig:vol06ch06:bonexs} shows a schematic mid-diaphyseal
cross-section. The cortical shell, $\sim 5\,\text{mm}$ thick in an
adult femur, carries most of the bending and torsional load through
the long bones; the trabecular interior, $\sim 70\,\%$ porous, carries
compression near the joint surfaces with the strut orientation
following the principal-stress trajectories. The marrow cavity is a
hydraulically incompressible filler with a haematopoietic function;
its compression resistance is engineering-negligible.

% Figure: schematic cross-section of a long bone (mid-diaphysis) showing
% the outer cortical (compact) shell, the inner trabecular (cancellous)
% region, and the marrow cavity. Wolff's-law alignment of trabecular
% struts indicated by the principal-stress trajectories.
\begin{figure}[htbp]
\centering
\begin{tikzpicture}[
  axis/.style={-{Stealth}, thick, draw=engblack},
  label/.style={font=\small, align=left},
]

% Outer cortical shell
\draw[thick, fill=engblack!10] (0, 0) circle (3.0);
% Inner cortical boundary
\draw[thick, fill=white] (0, 0) circle (2.4);
% Trabecular network region (outer ring of inner cavity)
\draw[thick, dashed, fill=engred!10] (0, 0) circle (1.7);
% Marrow cavity
\draw[thick, fill=englime!15] (0, 0) circle (0.9);

% Trabecular struts: a sparse pattern aligned with principal stresses
\foreach \a in {0, 30, 60, 90, 120, 150} {
  \draw[engred!70!black, thick] ({1.0*cos(\a)}, {1.0*sin(\a)})
    -- ({1.65*cos(\a)}, {1.65*sin(\a)});
  \draw[engred!70!black, thick] ({1.0*cos(\a+15)}, {1.0*sin(\a+15)})
    -- ({1.65*cos(\a+15)}, {1.65*sin(\a+15)});
}
% Transverse cross-struts
\foreach \r in {1.15, 1.40} {
  \draw[engred!50!black, thin] (0,0) circle (\r);
}

% Labels (each with align=left to permit \\ line breaks)
\node[anchor=west, align=left] at (3.4, 1.8) {Cortical bone\\(compact)};
\draw[-{Stealth}, thin] (3.3, 1.5) -- (2.7, 1.0);

\node[anchor=west, align=left] at (3.4, 0.0) {Trabecular bone\\(cancellous)};
\draw[-{Stealth}, thin] (3.3, -0.2) -- (1.6, -0.6);

\node[anchor=west, align=left] at (3.4, -1.6) {Marrow cavity};
\draw[-{Stealth}, thin] (3.3, -1.6) -- (0.7, -0.3);

% Scale bar
\draw[thick] (-3.2, -3.5) -- (-1.2, -3.5);
\node[anchor=north, font=\scriptsize] at (-2.2, -3.55) {$\sim 2\,\text{cm}$ (adult femur)};

% Properties annotations
\node[anchor=west, align=left, font=\scriptsize] at (-3.4, 3.7) {%
  $E_{\text{cort}} \approx 15\,\text{GPa}$, $\sigma_{\text{c}} \approx 200\,\text{MPa}$\\
  $E_{\text{trab}} \approx 0.5\text{ to }1.5\,\text{GPa}$\\
  $\rho_{\text{cort}} \approx 1.9\,\text{g/cm}^3$, $\rho_{\text{trab}} \approx 0.2\,\text{g/cm}^3$};

\end{tikzpicture}
\caption{Schematic mid-diaphyseal cross-section of an adult long bone.
The dense cortical shell carries most of the bending and torsional load;
the trabecular interior carries compression near the joint surfaces.
Trabecular struts align with the principal-stress trajectories
(Wolff's law). The marrow cavity is a hydraulically incompressible
filler with a haematopoietic function. Quantitative properties are
typical for the femoral mid-shaft of a $30$-year-old adult.}
\label{fig:vol06ch06:bonexs}
\end{figure}


\engterm{Cortical-bone stress-strain in detail}. Cortical bone is
anisotropic, viscoelastic, and stronger in compression than in
tension. Working numbers from the Reilly and Burstein 1975 study of
adult human femoral cortical bone
\cite[Tab.~3]{paper:v6c06-reilly-burstein-1975}:

\begin{center}
\begin{tabular}{lrrrr}
\toprule
Loading mode & $E$ (GPa) & UTS (MPa) & $\varepsilon_{\text{f}}$ & $\nu$ \\
\midrule
Longitudinal tension     & $17.0$ & $135$ & $0.030$ & $0.40$ \\
Longitudinal compression & $18.0$ & $205$ & $0.020$ & $0.40$ \\
Transverse tension       & $11.5$ & $\phantom{0}53$ & $0.020$ & $0.30$ \\
Transverse compression   & $12.0$ & $131$ & $0.025$ & $0.30$ \\
Shear (torsion)          & $\phantom{0}3.3$ & $\phantom{0}65$ & $0.027$ & n/a \\
\bottomrule
\end{tabular}
\end{center}

The compressive ultimate of $200\,\text{MPa}$ along the long axis is
the working figure for the chapter. The strain at compressive failure
is $\sim 0.02$, so cortical bone is intermediate between concrete
(failure strain $\sim 0.003$) and ductile steel (failure strain
$\sim 0.20$): brittle at the structural scale, but with five to ten
times more strain-to-failure than unreinforced ceramic. The tensile
strength is significantly lower than the compressive, reflecting the
typical composite-material asymmetry. Trabecular bone is approximately
fifty times more compliant than cortical bone at low apparent
density, with modulus scaling as
$E_{\text{trab}} \propto \rho^{2.5}$ over the physiological density
range \cite[Tab.~3.1]{paper:v6c06-hayes-1991}. The power-law form is
the key empirical fact for assessing osteoporotic fracture risk:
$10\,\%$ loss of trabecular bone density produces $\sim 25\,\%$ loss
of trabecular modulus and a corresponding drop in compressive
strength.

\engterm{Joints}. Synovial joints (knees, hips, shoulders, elbows,
fingers) consist of articular cartilage on the bony surfaces, a
synovial membrane that produces the lubricating synovial fluid, a
joint capsule that contains the fluid, and ligaments that constrain
the joint to its physiological motion. Articular cartilage is a
highly hydrated proteoglycan-collagen composite with Young's modulus
of $\sim 0.5$ to $1\,\text{MPa}$ and a coefficient of friction below
$0.01$ when properly lubricated. The lubrication is a combination of
boundary lubrication (hyaluronic acid and lubricin molecules adsorbed
to the cartilage surface) and hydrodynamic lubrication (the synovial
fluid film in motion). The combination is unusually low-friction; the
engineering of a synthetic joint replacement (Section~\ref{vol06:ch06}.7)
accepts substantially higher friction as a cost of the prosthesis.

\engterm{Muscles}. Skeletal muscle is the active actuator of the
musculoskeletal system. A muscle is composed of bundles of muscle
fibres (multinucleate muscle cells), each containing thousands of
myofibrils, each composed of sarcomeres in series
(Chapter~\ref{vol06:ch04}). Three principal fibre types occur: type
I (slow oxidative, fatigue-resistant, low maximum force), type IIa
(fast oxidative-glycolytic, intermediate), and type IIb (fast
glycolytic, high force, low endurance). Muscle types are mixed in
varying proportions across the body and across individuals; the
proportions adapt to training. Muscle force is generated by
ATP-driven actin-myosin crossbridges (Chapter~\ref{vol06:ch04});
maximum specific tension is approximately $0.3\,\text{MPa}$ in
human skeletal muscle \cite[p.~578]{paper:v6c06-bottinelli-1996}.

\engterm{The Hill model: force-length, force-velocity, activation}.
A working muscle force model takes three inputs (the contractile
element's length $L$, its shortening velocity $v$, and the activation
level $a$ between $0$ and $1$) and returns the muscle's force. The
Hill three-element model resolves the muscle into a contractile
element (CE) in parallel with a passive elastic element (PEE), in
series with a tendon series-elastic element (SEE). The active CE
force is the product of three factors:

\[
F_{\text{CE}}(L, v, a) = a \cdot F_0 \cdot f_L\!\left(\frac{L}{L_0}\right)
                       \cdot f_V\!\left(\frac{v}{v_{\max}}\right),
\]

where $F_0$ is the isometric peak force (the product of specific
tension $\sim 0.3\,\text{MPa}$ and physiological cross-section area
PCSA), $L_0$ the optimal length, $v_{\max}$ the maximum shortening
velocity ($2$ to $8\,L_0/\text{s}$ depending on fibre type), $f_L$ a
bell curve centred at $L/L_0 = 1$, and $f_V$ the Hill hyperbola
$(F + a^{*})(v + b^{*}) = (F_0 + a^{*}) b^{*}$ with the empirical
constants $a^{*}/F_0 \approx 0.25$ and $b^{*}/v_{\max} \approx 0.25$
\cite[p.~150]{paper:v6c06-hill-1938}. Eccentric (lengthening) muscle
force rises above $F_0$ to $\sim 1.5\,F_0$ to $1.8\,F_0$. The Hill
model is normalised in $F_0$, $L_0$, $v_{\max}$ and ports across
muscles with three numbers (PCSA, fibre length, fibre-type
fraction); Zajac's 1989 review gives the working parameter table
\cite[Tab.~II]{paper:v6c06-zajac-1989}.
Figure~\ref{fig:vol06ch06:muscle-curves} shows the two characteristic
curves; the script \texttt{hill\_muscle\_model.py} computes both and
applies them to a push-off estimate.

% Figure: muscle force-length and force-velocity curves (Hill model).
% Two panels side by side.
\begin{figure}[htbp]
\centering
\begin{tikzpicture}[
  axis/.style={-{Stealth}, thick, draw=engblack},
  curve/.style={very thick},
  annot/.style={font=\scriptsize, align=center},
]

% ===== Left panel: force-length =====
\begin{scope}[xshift=0cm]
\draw[axis] (0, 0) -- (6.2, 0) node[right, font=\small] {$L/L_0$};
\draw[axis] (0, 0) -- (0, 4.7) node[above, font=\small] {$F/F_0$};

\foreach \x/\xt in {1/0.5, 2/0.75, 3/1.0, 4/1.25, 5/1.5} {
  \draw (\x, 0) -- (\x, -0.08) node[below, font=\scriptsize] {\xt};
}
\foreach \y/\yt in {1/0.25, 2/0.5, 3/0.75, 4/1.0} {
  \draw (0, \y) -- (-0.08, \y) node[left, font=\scriptsize] {\yt};
}

% Active force-length: triangular about L/L0 = 1 with plateau
\draw[curve, engred, smooth] (1.0, 0) .. controls (2.0, 1.5) and (2.7, 3.5) .. (3.0, 4.0)
   -- (3.4, 4.0) .. controls (3.7, 4.0) and (4.5, 1.0) .. (5.0, 0);
\node[engred, font=\small, anchor=south] at (3.2, 4.0) {active};

% Passive force-length: rising exponentially above L/L0 ~ 1
\draw[curve, engblue, smooth, samples=50, domain=2.7:5.2]
  plot (\x, {0.05 * exp(1.0 * (\x - 2.7))});
\node[engblue, font=\small, anchor=west] at (5.0, 3.2) {passive};

% Total = active + passive (dashed)
\draw[curve, engblack, dashed, smooth, samples=50, domain=1.0:5.2]
  plot (\x, { (\x < 3.0 ? \x - 1.0 : (\x < 3.4 ? 4.0 : 4.0 - 8.0*(\x - 3.4)/(5.0 - 3.4))) + 0.05*exp(1.0*max(\x-2.7,0))});

\node[font=\bfseries, anchor=north] at (3.0, -0.7) {Force-length};
\end{scope}

% ===== Right panel: force-velocity =====
\begin{scope}[xshift=8.5cm]
\draw[axis] (0, 0) -- (6.2, 0) node[right, font=\small] {$v / v_{\max}$};
\draw[axis] (0, 0) -- (0, 4.7) node[above, font=\small] {$F / F_0$};

\foreach \x/\xt in {1/-0.5, 2/0, 3/0.5, 4/1.0} {
  \draw (\x+1.0, 0) -- (\x+1.0, -0.08) node[below, font=\scriptsize] {\xt};
}
\foreach \y/\yt in {1/0.5, 2/1.0, 3/1.5, 4/1.8} {
  \draw (0, \y) -- (-0.08, \y) node[left, font=\scriptsize] {\yt};
}

% Hill curve: concentric (v > 0): (F + a)(v + b) = (F_0 + a) b
% Take a/F0 = 0.25, b/vmax = 0.25
% F/F0 = ((F0+a) b / (v+b) - a) / F0
% Eccentric (v < 0): rises to 1.5-1.8 F0 then plateau
\draw[curve, engred, smooth, samples=80, domain=0:1.0]
  plot ({2.0 + 4.0*\x}, {2.0 * (1.25 * 0.25 / (\x + 0.25) - 0.25)});

% Eccentric branch
\draw[curve, engred, smooth, samples=40, domain=-0.5:0]
  plot ({2.0 + 4.0*\x}, {2.0 + 1.6 * (1 - exp(5*\x))});

% Mark isometric F_0 at v = 0
\fill[engred] (2.0, 2.0) circle (1.5pt);
\node[engred, anchor=south west, font=\scriptsize] at (2.1, 2.05) {$F_0$};
% Mark v_max at F = 0
\fill[engred] (6.0, 0) circle (1.5pt);
\node[engred, anchor=south, font=\scriptsize] at (6.0, 0.1) {$v_{\max}$};

\node[engred, font=\small, anchor=west] at (4.0, 0.9) {concentric};
\node[engred, font=\small, anchor=east] at (1.7, 3.4) {eccentric};

\node[font=\bfseries, anchor=north] at (3.0, -0.7) {Force-velocity (Hill)};
\end{scope}

\end{tikzpicture}
\caption{Hill-type muscle model: force-length (left) and force-velocity
(right) characteristics. Active force peaks at the optimal length
$L_0$ where actin-myosin overlap is maximal; passive force grows
exponentially above resting length as the connective-tissue elements
stretch. Concentric shortening at velocity $v$ produces force
$F(v) = (F_0 + a)\,b / (v + b) - a$, the Hill hyperbola, with
$a/F_0 \approx 0.25$ and $b/v_{\max} \approx 0.25$ for mammalian
skeletal muscle. Eccentric (lengthening) contraction produces forces
up to $\sim 1.5\,F_0$ to $1.8\,F_0$, the mechanical mechanism by
which downhill walking and the deceleration phase of running impose
the largest forces on tendons.}
\label{fig:vol06ch06:muscle-curves}
\end{figure}


\engterm{Tendons and ligaments}. Tendons connect muscle to bone;
ligaments connect bone to bone. Both are dense collagenous tissues
with collagen fibres oriented along the load direction. Tendons have
Young's modulus of $\sim 1\,\text{GPa}$ along the fibre, ultimate
tensile strength of $\sim 100\,\text{MPa}$, and stretch by $4\,\%$
to $8\,\%$ at physiological loads. The strain-energy storage in
tendons is functionally important: the Achilles tendon stores
$\sim 35\,\text{J}$ per stride during running and returns
$\sim 30\,\%$ of that energy on the next stride, reducing the
metabolic cost of running by approximately $35\,\%$ compared to a
hypothetical inelastic equivalent
\cite[ch.~11]{text:winter-biomechanics-human},
\cite[p.~62]{paper:v6c06-alexander-1991}.

\engterm{The skeleton as a structural-mechanics problem}. The body's
skeleton can be analysed by the methods of Volume III. Bones are
loaded primarily in compression at the joints and by combined
compression-bending in the long bones; joints are subject to combined
compression and shear; tendons and ligaments to tension; muscles to
tension only. The free-body diagram of a body segment is the working
tool for analysing musculoskeletal forces. A typical inverse-dynamics
analysis of human walking (the gait-analysis pipeline that uses
ground-reaction-force plates, motion capture, and segmental dynamics)
produces joint moments at every $1\,\text{ms}$ during the gait cycle
and constitutes the standard biomechanical assessment in orthopaedics
and rehabilitation.

\begin{mastery}
\textit{The lever-arm trade-off in vertebrate limbs.} The skeletal
geometry of every vertebrate limb makes the same engineering trade:
the muscle-tendon insertion sits close to the joint centre
(small moment arm, $\sim 5\,\text{cm}$ at the patellar tendon), and
the limb extends far from the joint (large moment arm,
$\sim 40\,\text{cm}$ at the foot from the knee centre). The ratio of
$8{:}1$ means the muscle must produce eight times the external force
in tendon tension to balance the external moment. The cost is high
internal force (and the corresponding high joint compression that
drives the patello-femoral joint pressure to $\sim 22\,\text{MPa}$ in
deep flexion). The benefit is an $8{:}1$ velocity ratio in the
opposite direction: a $5\,\text{cm/s}$ muscle contraction produces a
$40\,\text{cm/s}$ foot velocity. The trade is the same one a robotic
arm with a strong, slow motor and a long output lever would make,
and is the geometric reason that animals are designed for high tip
velocity at the cost of high joint loads.
\end{mastery}

\section{Locomotion: gait, swimming, flight}\pathtag{standard}

Animals move by exerting forces on the environment. The three
principal modes (legged walking and running, swimming, flying) each
present a distinct engineering problem and a distinct set of
trade-offs.

\engterm{Human gait}. Human walking is an inverted-pendulum motion:
the body's centre of mass arcs upward and forward over each stance
leg, falling under gravity into the next step. The cost is the
gravitational potential energy lost at each stance phase; the
gain is the slow, low-impact, energetically efficient locomotion
that walking provides. The metabolic cost of walking is approximately
$0.5\,\text{J}/(\text{kg}\cdot\text{m})$ at the optimum walking
speed of $1.4\,\text{m}/\text{s}$ for a typical adult. The transition
from walking to running occurs near $2.0\,\text{m}/\text{s}$ and is
marked by a change in mechanics from inverted-pendulum (out-of-phase
kinetic and potential energy) to mass-spring (in-phase kinetic and
potential energy, with the elastic energy of the leg tendons storing
and returning energy in each step)
\cite[p.~R250]{paper:v6c06-cavagna-1977}. Running has higher metabolic
cost per metre at low speeds but the cost is approximately constant
with speed in running, so above the transition speed running becomes
more efficient than walking on a per-time basis
\cite[ch.~9]{text:winter-biomechanics-human}.

Figure~\ref{fig:vol06ch06:gait-phases} schematises the gait phases
and the corresponding vertical ground-reaction force trace. The
double-hump GRF (first peak at loading response, valley at mid-stance,
second peak at push-off) is the signature of normal walking and the
first thing a clinical gait analyst inspects. In running, the GRF
collapses to a single peak of $2.5$ to $3.5\,\text{BW}$; the absence
of the double valley is the mechanical signature of the mass-spring
gait.

% Figure: phases of the gait cycle, heel strike to toe off, with the
% percent-cycle annotation and the ground-reaction-force trace.
\begin{figure}[htbp]
\centering
\begin{tikzpicture}[
  axis/.style={-{Stealth}, thick, draw=engblack},
  stick/.style={thick, draw=engblack},
  label/.style={font=\scriptsize, align=center},
]

% Five stick-figure poses across the page representing gait phases.
\foreach \i/\x/\hipy/\kneey/\footx/\caption in {%
  0/0/3.4/1.7/0.2/heel strike,
  1/3.0/3.2/1.5/-0.1/loading response,
  2/6.0/3.0/1.4/-0.2/mid stance,
  3/9.0/3.2/1.5/-0.4/terminal stance,
  4/12.0/3.4/1.7/-0.6/toe off}
{
  % Head
  \draw[stick, fill=engblack!10] (\x, 4.2) circle (0.25);
  % Trunk
  \draw[stick] (\x, 3.95) -- (\x, \hipy);
  % Pelvis
  \draw[stick] (\x-0.2, \hipy) -- (\x+0.2, \hipy);
  % Stance leg
  \draw[stick, engred, very thick] (\x-0.2, \hipy) -- (\x+\footx-0.15, \kneey) -- (\x+\footx, 0.4);
  % Foot
  \draw[stick, engred, very thick] (\x+\footx-0.2, 0.4) -- (\x+\footx+0.35, 0.4);
  % Swing leg (mirror, lower stance for first/last)
  \draw[stick] (\x+0.2, \hipy) -- (\x-\footx+0.1, \kneey - 0.3) -- (\x-\footx+0.05, 0.55);
  % Caption
  \node[label] at (\x, -0.2) {\caption\\$\i\times25\,\%$};
}

% GRF curve underneath (two humps, normalised)
\begin{scope}[yshift=-3.5cm]
\draw[axis] (-0.5, 0) -- (13.5, 0) node[right, font=\small] {$t$ (\% gait cycle)};
\draw[axis] (-0.5, 0) -- (-0.5, 3.0) node[above, font=\small] {GRF $/$ BW};

\foreach \x/\xt in {0/0, 3/25, 6/50, 9/75, 12/100} {
  \draw (\x, 0) -- (\x, -0.08) node[below, font=\scriptsize] {\xt};
}
\foreach \y/\yt in {1/0.5, 2/1.0, 2.4/1.2} {
  \draw (-0.5, \y) -- (-0.6, \y) node[left, font=\scriptsize] {\yt};
}

% Double-hump curve typical of walking GRF
\draw[very thick, engblue, smooth, samples=100, domain=0:12]
  plot (\x, {2.4 * exp(-((\x - 2.0)/1.2)^2) + 2.2 * exp(-((\x - 8.0)/1.2)^2)});

\node[anchor=west, font=\small, engblue] at (2.0, 2.7) {first peak};
\node[anchor=west, font=\small, engblue] at (8.0, 2.5) {push-off};
\node[anchor=west, font=\small] at (5.0, 1.0) {valley};
\end{scope}

\end{tikzpicture}
\caption{Phases of the gait cycle (single stance leg, red) and the
typical vertical ground-reaction force trace below. The double-hump
GRF is the signature of normal walking: the first peak occurs at
loading response (the body's centre of mass falls onto the stance
leg), the valley at mid-stance (the centre of mass rises over the
stance leg, momentarily decelerating the vertical fall), and the
second peak at push-off (the ankle plantar-flexors propel the body
forward and up). Peak vertical GRF is typically $1.1$ to $1.3\,\text{BW}$
in walking; in running, the curve collapses to a single peak of
$2.5$ to $3.5\,\text{BW}$. The gait cycle conventionally runs from
heel strike (or initial contact) of one foot to the next heel strike
of the same foot.}
\label{fig:vol06ch06:gait-phases}
\end{figure}


\engterm{Gait analysis}. The standard biomechanical assessment of
human gait combines optical motion capture (tracking markers on the
body segments), force plates (measuring ground reaction forces), and
electromyography (measuring muscle activation patterns). The
recorded data are processed by inverse dynamics to compute joint
moments and powers throughout the gait cycle. The result is a
quantitative summary of the patient's gait that supports orthotic
design, prosthetic fitting, and surgical planning. Modern gait
analysis is the standard pre-operative assessment for paediatric
cerebral palsy, post-stroke rehabilitation, and lower-limb
amputation.

\engterm{Swimming at large and small scales}. Swimming biomechanics
differ across the Reynolds-number range. At large Reynolds number
($\text{Re} > 10^3$, fish and aquatic mammals), inertia dominates
viscosity, and the body's wake carries momentum that the swimmer
uses to generate thrust. At small Reynolds number
($\text{Re} < 1$, bacteria and protozoa), viscosity dominates and the
swimmer must use non-reciprocal motion to break time-reversal
symmetry; a single oscillating flagellum is the canonical bacterial
solution. The reader who wishes to understand the small-Reynolds
regime should consult Purcell's 1977 paper ``Life at low Reynolds
number'' \cite[p.~5]{paper:v6c06-purcell-1977}, which lays out the
scallop theorem and the engineering of bacterial swimming.

\engterm{Animal flight}. Flying animals span four orders of magnitude
in size, from $1\,\text{mg}$ insects to $15\,\text{kg}$ albatrosses.
The wing's lift force scales with wing area and air-speed squared;
the body's weight scales with body volume (mass). For an animal that
must support its weight in flight, the wing-loading (mass per unit
wing area) and the flight speed scale predictably with body size:
larger animals fly faster and have higher wing-loading. The
allometric relations match the predictions of basic aerodynamics
across the size range \cite[ch.~6]{text:fung-biomechanics}.

\section{Cardiovascular biomechanics}\pathtag{standard}

The cardiovascular system is a pulsatile hydraulic network that
delivers oxygen and nutrients to all body tissues and removes carbon
dioxide and metabolic waste. The system has two principal pumps (the
heart's left and right ventricles), a high-pressure distribution
network (the arterial tree), a high-surface-area exchange network
(the capillary beds), and a low-pressure return network (the venous
tree). The engineering of the system is one of the cleanest
applications of fluid mechanics to a biological problem
\cite[ch.~5]{text:fung-biomechanics},
\cite[ch.~3]{text:v6c06-klabunde-cv-2011}.

\engterm{The heart as a pump}. The left ventricle pumps blood at a
mean pressure of $\sim 100\,\text{mmHg}$ (systolic $120\,\text{mmHg}$,
diastolic $80\,\text{mmHg}$) and at a flow rate of $\sim 5\,\text{L}/
\text{min}$ at rest, rising to $25\,\text{L}/\text{min}$ during heavy
exercise. The stroke volume is $\sim 70\,\text{mL}$; the heart rate
is $60$ to $80\,\text{beats}/\text{min}$ at rest. The cardiac power
output is $\sim 1$ to $2\,\text{W}$ at rest, $\sim 10\,\text{W}$
during heavy exercise. The pump is the body's most metabolically
active organ per unit mass: at $\sim 300\,\text{g}$, the heart
consumes $\sim 7\,\%$ of total body oxygen.

\engterm{Pressure-volume loop}. The heart's behaviour is summarised by
the pressure-volume loop, which plots ventricular pressure against
ventricular volume over one heartbeat. The loop has four phases:
diastolic filling (low pressure, increasing volume), isovolumetric
contraction (rising pressure at constant volume), ejection (high
pressure, decreasing volume), and isovolumetric relaxation (falling
pressure at constant volume). The area inside the loop is the
stroke work performed by the ventricle. Two reference curves frame
the loop: the end-diastolic pressure-volume relation (EDPVR, the
passive filling curve, governed by the ventricular wall's
passive stiffness), and the end-systolic pressure-volume relation
(ESPVR, a near-linear curve whose slope $E_{\text{es}}$ is the
load-independent measure of contractility introduced by Suga and
Sagawa in $1974$
\cite[p.~122]{paper:v6c06-suga-sagawa-1974}). The loop's geometry is
the clinical fingerprint of cardiac function: dilated
cardiomyopathy shifts the loop rightward (larger volumes); concentric
hypertrophy shifts the loop leftward and upward; failing hearts
compress the loop's stroke volume.
Figure~\ref{fig:vol06ch06:pv-loop} shows the four phases and the
two reference curves.

% Figure: Left-ventricular pressure-volume loop with the four phases
% (filling, isovolumetric contraction, ejection, isovolumetric
% relaxation) annotated.
\begin{figure}[htbp]
\centering
\begin{tikzpicture}[
  axis/.style={-{Stealth}, thick, draw=engblack},
  loop/.style={very thick, engred},
  edge/.style={thin, engblack},
  annot/.style={font=\scriptsize, align=left},
]

\draw[axis] (0, 0) -- (10, 0) node[right, font=\small] {LV volume (mL)};
\draw[axis] (0, 0) -- (0, 7) node[above, font=\small] {LV pressure (mmHg)};

\foreach \x/\xt in {2/40, 4/80, 6/120, 8/160} {
  \draw (\x, 0) -- (\x, -0.1) node[below, font=\scriptsize] {\xt};
}
\foreach \y/\yt in {1/20, 2.5/50, 4/80, 5.5/110, 6.5/140} {
  \draw (0, \y) -- (-0.1, \y) node[left, font=\scriptsize] {\yt};
}

% PV loop: simple parallelogram-with-curves approximation.
% Vertices: A = filling start (50 mL, 5 mmHg)
%           B = end-diastole / mitral close (120 mL, 8 mmHg)
%           C = aortic open (120 mL, 80 mmHg)
%           D = end-systole / aortic close (50 mL, 100 mmHg)
% A: (2.5, 0.25), B: (6, 0.4), C: (6, 4), D: (2.5, 5)

% Filling (A->B): along EDPVR (passive end-diastolic curve)
\draw[loop] (2.5, 0.25) .. controls (4.5, 0.3) and (5.5, 0.35) .. (6, 0.4);
% Isovolumetric contraction (B->C)
\draw[loop] (6, 0.4) -- (6, 4);
% Ejection (C->D): along the high-pressure plateau, slight curve
\draw[loop] (6, 4) .. controls (5.0, 4.8) and (3.5, 5.1) .. (2.5, 5);
% Isovolumetric relaxation (D->A)
\draw[loop] (2.5, 5) -- (2.5, 0.25);

% ESPVR (end-systolic pressure-volume relation): linear, dashed
\draw[edge, dashed] (1.0, 0) -- (4.0, 7);
\node[annot, anchor=west] at (3.5, 6.5) {ESPVR};

% EDPVR: low and nearly flat
\draw[edge, dashed] (1.0, 0) .. controls (4.0, 0.1) and (6.5, 0.5) .. (8.0, 1.7);
\node[annot, anchor=west] at (7.5, 1.3) {EDPVR};

% Vertex labels
\fill[engred] (2.5, 0.25) circle (2pt);
\node[annot, anchor=east] at (2.4, 0.4) {A (MV open)};
\fill[engred] (6, 0.4) circle (2pt);
\node[annot, anchor=west] at (6.1, 0.3) {B (MV close)};
\fill[engred] (6, 4) circle (2pt);
\node[annot, anchor=west] at (6.1, 4.0) {C (AV open)};
\fill[engred] (2.5, 5) circle (2pt);
\node[annot, anchor=east] at (2.4, 5.2) {D (AV close)};

% Stroke volume and stroke work
\node[font=\small, engred] at (4.3, 2.7) {SV};
\node[annot] at (4.3, 2.3) {stroke work\\$\approx 1\,\text{J}/\text{beat}$};

% Phase annotations
\node[annot] at (7.5, 0.1) {filling (A$\to$B)};
\node[annot] at (6.5, 2.5) {iso. contr.\\(B$\to$C)};
\node[annot] at (4.3, 4.8) {ejection (C$\to$D)};
\node[annot] at (1.0, 2.5) {iso.\\relax.\\(D$\to$A)};

\node[font=\bfseries, anchor=south] at (5, 7.0) {Left ventricular pressure-volume loop};

\end{tikzpicture}
\caption{Left-ventricular pressure-volume loop, one cardiac cycle.
Phase A$\to$B is diastolic filling, governed by the end-diastolic
pressure-volume relation (EDPVR, lower dashed curve). Phase B$\to$C
is isovolumetric contraction (mitral valve closed, aortic valve not
yet open). Phase C$\to$D is ejection. Phase D$\to$A is isovolumetric
relaxation. End-systolic points lie on the end-systolic
pressure-volume relation (ESPVR, upper dashed line), whose slope is
the load-independent measure of contractility. The area inside the
loop is the stroke work per beat, $\approx 1\,\text{J}$ at rest. The
loop's geometry is the clinical fingerprint of cardiac function:
dilation moves it rightward, concentric hypertrophy upward and
leftward, failing hearts collapse the stroke volume.}
\label{fig:vol06ch06:pv-loop}
\end{figure}


The script \texttt{pressure\_volume\_loop.py} computes the standard
loop summary (SV, EF, stroke work, cardiac output, cardiac power) at
rest, during heavy exercise, and in heart failure with reduced
ejection fraction. Sample output: at rest, $SV = 70\,\text{mL}$,
$EF = 58\,\%$, $W_s \approx 0.93\,\text{J/beat}$, cardiac power
$\approx 1.1\,\text{W}$. At heavy exercise (with sympathetic
contractility increase), $SV = 105\,\text{mL}$, $EF = 72\,\%$,
$W_s \approx 1.85\,\text{J/beat}$, cardiac power
$\approx 5.2\,\text{W}$. In heart failure, $SV = 60\,\text{mL}$,
$EF = 33\,\%$, $W_s \approx 0.68\,\text{J/beat}$ at a higher EDV
($180\,\text{mL}$): the loop has shifted rightward with reduced
stroke volume.

\engterm{Murray's law and branching geometry}. The arterial tree
branches with a regular geometric rule. Murray's $1926$ argument
minimises the sum of the metabolic cost of maintaining the blood
volume (proportional to lumen volume, $\sim d^2 L$) and the
Poiseuille pumping work (inversely proportional to $d^4$), and finds
that the diameters at a bifurcation satisfy
\[
d_{\text{parent}}^3 = d_1^3 + d_2^3,
\]
so a symmetric daughter has $d_{\text{d}}/d_{\text{p}} =
2^{-1/3} \approx 0.794$
\cite[p.~209]{paper:v6c06-murray-1926}. The corresponding
shear-stress rule is that wall shear stress is preserved across
generations (the optimal balance is also the constant-shear
balance). The optimal symmetric bifurcation half-angle satisfies
$\cos(2\theta) = 2^{1/3} - 1$, or $\theta \approx 37.5^{\circ}$.
Figure~\ref{fig:vol06ch06:cv-tree} schematises the hierarchy and the
total cross-sectional area / velocity inversion across the bed.

% Figure: schematic cardiovascular branching tree with Murray's-law
% radius scaling annotations and cross-sectional area / velocity inset.
\begin{figure}[htbp]
\centering
\begin{tikzpicture}[
  axis/.style={-{Stealth}, thick, draw=engblack},
  vessel/.style={engred, line cap=round},
  label/.style={font=\scriptsize, align=left},
]

% Vessel hierarchy as horizontal segments with radius proportional to
% the diameter scale.
\foreach \i/\name/\d/\v/\area in {%
  0/aorta/0.50/0.5/4.5,
  1/large artery/0.30/0.35/15,
  2/arteriole/0.15/0.05/400,
  3/capillary/0.05/0.001/4500,
  4/venule/0.10/0.005/2500,
  5/large vein/0.30/0.15/40,
  6/vena cava/0.45/0.3/10}
{
  \pgfmathsetmacro{\xv}{2.0 * \i}
  \draw[vessel, line width=\d cm] (\xv, 3.0) -- (\xv+1.8, 3.0);
  \node[label] at (\xv+0.9, 3.6) {\name};
  \node[label, anchor=north] at (\xv+0.9, 2.85)
    {$d\sim$ \pgfmathprintnumber[fixed,precision=2]{\d}\,cm};
}

% Branching tree from aorta -> arteries -> arterioles -> capillaries
% Below the labelled hierarchy.
\begin{scope}[yshift=-1.3cm]
\node[engred, font=\bfseries] at (-0.5, 1.5) {Heart};
\draw[engred, very thick] (0, 1.5) -- (2, 1.5);

% Level 1 (large artery)
\foreach \i in {0, 1} {
  \draw[engred, thick] (2, 1.5) -- (3.5, 1.5 + 0.7*(\i - 0.5));
}
% Level 2 (arteriole)
\foreach \i in {0, 1, 2, 3} {
  \draw[engred!80, thick] (3.5, 1.85) -- (5.0, 1.85 + 0.3*(\i - 1.5));
  \draw[engred!80, thick] (3.5, 1.15) -- (5.0, 1.15 + 0.3*(\i - 1.5));
}
% Level 3 (capillary, drawn as dots beyond)
\foreach \y in {0.4, 0.7, 1.0, 1.3, 1.6, 1.9, 2.2, 2.5} {
  \draw[engred!50, thin] (5.0, \y) -- (6.5, \y);
}

\node[label, anchor=west] at (6.7, 1.5)
  {$d_{\text{p}}^3 = d_1^3 + d_2^3$\\Murray's law\\$d_{\text{daughter}}/d_{\text{parent}} = 2^{-1/3} \approx 0.794$};
\end{scope}

% Inset: cross-sectional area vs velocity
\begin{scope}[xshift=0cm, yshift=-5.5cm, scale=0.7]
\draw[axis] (0, 0) -- (10, 0) node[right, font=\small] {vessel generation};
\draw[axis] (0, 0) -- (0, 4) node[above, font=\small] {(log scale)};

% Plot total cross-sectional area (rising to capillaries, falling on return)
\draw[very thick, engblue, smooth] (0, 0.5) .. controls (1.5, 0.8) and (3.5, 2.2)
    .. (4.5, 3.5) .. controls (5.5, 2.5) and (7, 1.2) .. (9, 0.6);
\node[engblue, font=\small, anchor=south] at (4.5, 3.55) {total area};

% Velocity (falling, then rising)
\draw[very thick, engred, dashed, smooth] (0, 3.5) .. controls (1.5, 3.0) and (3.5, 1.0)
    .. (4.5, 0.5) .. controls (5.5, 1.0) and (7, 2.0) .. (9, 2.5);
\node[engred, font=\small, anchor=west] at (9.0, 2.55) {velocity};

\foreach \x/\xt in {0/aorta, 2/artery, 4/capillary, 6/venule, 8/vein} {
  \draw (\x, 0) -- (\x, -0.1) node[below, font=\tiny] {\xt};
}
\end{scope}

\end{tikzpicture}
\caption{Cardiovascular branching tree. Top: vessel hierarchy with
characteristic diameters from aorta ($\sim 1\,\text{cm}$) to capillary
($\sim 8\,\si{\micro\meter}$) and back through venules to the vena
cava. The diameter ratio at each symmetric bifurcation follows
Murray's law (the cubed-radius constraint that minimises the sum of
metabolic and pumping cost). Bottom inset: total cross-sectional area
rises by three orders of magnitude from aorta to capillary bed, and
blood velocity falls by the same factor (mass conservation in an
incompressible fluid). Slow capillary transit, $\sim 1\,\text{s}$, is
the time scale the exchange biochemistry can use.}
\label{fig:vol06ch06:cv-tree}
\end{figure}


The script \texttt{murray\_branching.py} tabulates the diameter,
total cross-sectional area, blood velocity, and Reynolds number
across the $24$-generation symmetric tree from the aorta to the
capillaries. The Reynolds number falls from $\text{Re} \approx 4000$
at the aorta ($d = 25\,\text{mm}$, $v \approx 17\,\text{cm/s}$) to
$\text{Re} \approx 0.001$ at the capillary
($d = 8\,\si{\micro\meter}$, $v \approx 0.05\,\text{cm/s}$). The
implication for capillary mechanics: blood flow at capillary scale
is creeping, and the bending and folding of erythrocytes through
capillaries narrower than the cell diameter is the dominant
microvascular flow phenomenon.

\engterm{Reynolds number in blood flow}. With $\rho = 1060\,\text{kg/m}^3$
and $\mu \approx 4\,\text{cP} = 4 \times 10^{-3}\,\text{Pa}\cdot\text{s}$
(whole-blood apparent viscosity at large shear rates):
\[
\text{Re}_{\text{capillary}} = \frac{\rho v d}{\mu}
= \frac{1060 \cdot 5\times 10^{-4} \cdot 8\times 10^{-6}}{4\times 10^{-3}}
\approx 1 \times 10^{-3}.
\]
The capillary flow is in the low-Reynolds regime (Stokes flow);
inertia is negligible, and the cell-membrane mechanics of the
erythrocyte squeezing through the $4$-$\si{\micro\meter}$ lumen
dominate the apparent viscosity. The Fahraeus-Lindqvist effect
(apparent viscosity falling with tube radius below
$\sim 300\,\si{\micro\meter}$) is the consequence: erythrocytes
align centrally and leave a cell-poor plasma layer near the wall.

\engterm{Arterial wave propagation and the Moens-Korteweg equation}.
Blood pressure propagates through the arterial tree as a wave. The
pressure wave's velocity is the pulse-wave velocity (PWV), typically
$5$ to $10\,\text{m}/\text{s}$ in healthy adults. PWV depends on the
artery's stiffness through the Moens-Korteweg equation
\cite[p.~538]{paper:v6c06-korteweg-1878},
\cite[p.~24]{paper:v6c06-moens-1878}
\[
c = \sqrt{\frac{E h}{2 \rho R}},
\]
where $E$ is the artery's incremental elastic modulus, $h$ is the
wall thickness, $\rho$ is the blood density, and $R$ is the artery's
inner radius. The derivation balances pressure-driven wall
distension against fluid mass conservation in a thin elastic tube;
the assumptions are thin wall ($h \ll R$), linear elasticity,
incompressible fluid, no viscous losses, and small amplitude. PWV
rises with age (from $\sim 5\,\text{m}/\text{s}$ in adolescents to
$\sim 10\,\text{m}/\text{s}$ in seventy-year-olds) because the
arteries stiffen progressively. PWV is a clinical biomarker of
arterial stiffening and a strong predictor of cardiovascular event
risk.

\engterm{Capillary exchange}. The capillary network is the
exchange interface between blood and tissue. Capillaries are
$8\,\si{\micro\meter}$ in diameter, only slightly larger than a red
blood cell, and have walls of a single layer of endothelial cells.
The combined cross-sectional area of all capillaries
($\sim 4500\,\text{cm}^2$ in the body) is approximately $1000$ times
the cross-section of the aorta, so blood velocity falls dramatically
in the capillaries (from $\sim 50\,\text{cm}/\text{s}$ in the aorta
to $\sim 0.05\,\text{cm}/\text{s}$ in capillaries). The slower
velocity is essential: capillary transit times are $\sim 1\,\text{s}$,
which is the time needed for oxygen exchange between the red blood
cell and the surrounding tissue.

\engterm{Starling forces}. Fluid exchange across the capillary wall
is set by the balance between hydrostatic and osmotic pressure
differences (the Starling forces). Hydrostatic pressure pushes fluid
out of the capillary; the protein-based osmotic pressure (colloid
osmotic pressure, $\sim 25\,\text{mmHg}$) pulls it back. The
net filtration over the body's capillary surface is
$\sim 2\,\text{L}/\text{day}$, balanced by lymphatic return. Failure
of the balance produces oedema (excess interstitial fluid), a
common clinical sign in heart failure, kidney failure, and
inflammatory states.

\section{Respiratory biomechanics}\pathtag{standard}

The respiratory system exchanges oxygen and carbon dioxide between
the blood and the atmosphere. The engineering of the system is
constrained by the diffusion-limited nature of gas exchange and by
the mechanical work required to ventilate the lungs.

\engterm{Anatomy and scaling}. The human lung contains
$\sim 3 \times 10^8$ alveoli with a total exchange surface area of
$70$ to $100\,\text{m}^2$. The alveoli are arranged in
$\sim 23$ generations of bifurcating airways from the trachea to
the terminal alveolar sacs. The scaling is dictated by the trade-off
between exchange surface area (scales with alveolar count) and
diffusion path length (which must be short for adequate exchange).

\engterm{Ventilation mechanics}. Inspiration is active: the
diaphragm contracts downward and the intercostal muscles raise the
ribcage, expanding the thoracic volume and reducing intra-thoracic
pressure below atmospheric. Air flows into the lungs along the
pressure gradient. Expiration is mostly passive: the elastic recoil
of the lungs and chest wall reduces the volume back to functional
residual capacity. The work of breathing is $\sim 1$ to $2\,\text{J}/
\text{breath}$ at rest, with a respiratory rate of $\sim 12\,
\text{breaths}/\text{min}$, totalling $\sim 1\,\%$ of resting
metabolic rate. The work rises rapidly during exercise and can
exceed $10\,\%$ of total metabolic rate during heavy work
\cite[ch.~7]{text:v6c06-west-respiratory-2012}.

\engterm{Compliance and resistance}. The lung's mechanical behaviour
is summarised by two parameters: \engterm{compliance}
$C = \Delta V/\Delta P \approx 200\,\text{mL}/\text{cmH}_2\text{O}$
(the volume change per unit pressure change), and \engterm{airway
resistance} $R \approx 1\,\text{cmH}_2\text{O}/(\text{L}/\text{s})$
(the pressure-flow relationship of the airways). The product $RC$
is the lung's time constant, $\sim 0.2\,\text{s}$, well below the
respiratory cycle period. Pathological changes in $C$ (emphysema
increases it; fibrosis decreases it) and $R$ (asthma and chronic
bronchitis increase it) are the principal mechanisms of obstructive
and restrictive pulmonary disease.
Figure~\ref{fig:vol06ch06:resp-mech} shows the static P-V relation
and the flow-volume loop in spirometry.

% Figure: Respiratory mechanics. Two panels: P-V curve of the lung
% (compliance) and flow-volume loop (a standard spirometry trace).
\begin{figure}[htbp]
\centering
\begin{tikzpicture}[
  axis/.style={-{Stealth}, thick, draw=engblack},
  curve/.style={very thick},
  annot/.style={font=\scriptsize, align=left},
]

% ===== Left panel: lung pressure-volume curve =====
\begin{scope}[xshift=0cm]
\draw[axis] (0, 0) -- (5.5, 0) node[right, font=\small] {$P_{\text{tp}}$ (cmH$_2$O)};
\draw[axis] (0, 0) -- (0, 5.5) node[above, font=\small] {$V$ (L)};

\foreach \x/\xt in {1/5, 2/10, 3/15, 4/20, 5/25} {
  \draw (\x, 0) -- (\x, -0.08) node[below, font=\scriptsize] {\xt};
}
\foreach \y/\yt in {1/1, 2/2, 3/3, 4/4, 5/5} {
  \draw (0, \y) -- (-0.08, \y) node[left, font=\scriptsize] {\yt};
}

% Inflation curve (lower): sigmoidal
\draw[curve, engred, smooth, samples=80, domain=0:5.0]
  plot (\x, {5.0 / (1.0 + exp(-1.6*(\x - 2.0)))});
% Deflation curve (upper) with hysteresis
\draw[curve, engblue, smooth, samples=80, domain=0:5.0]
  plot (\x, {5.0 / (1.0 + exp(-1.6*(\x - 1.5)))});

% Compliance triangle near FRC
\draw[engblack, dashed] (1.6, 1.5) -- (2.6, 1.5) -- (2.6, 3.0) -- cycle;
\node[annot] at (2.7, 2.2) {slope $= C \approx 0.2\,\frac{\text{L}}{\text{cmH}_2\text{O}}$};

\node[engred, font=\small, anchor=west] at (3.5, 3.2) {inflation};
\node[engblue, font=\small, anchor=west] at (1.5, 4.4) {deflation};

\node[font=\bfseries, anchor=south] at (2.7, 5.6) {Pressure-volume};
\node[font=\scriptsize, anchor=north] at (2.7, -0.7) {hysteresis $=$ work of breathing};
\end{scope}

% ===== Right panel: flow-volume loop (spirometry) =====
\begin{scope}[xshift=8cm]
\draw[axis] (0, 0) -- (5.5, 0) node[right, font=\small] {Volume (L)};
\draw[axis] (0, -3) -- (0, 4.0) node[above, font=\small] {Flow (L/s)};

\foreach \x/\xt in {1/1, 2/2, 3/3, 4/4, 5/5} {
  \draw (\x, 0) -- (\x, -0.08) node[below, font=\scriptsize] {\xt};
}
\foreach \y/\yt in {-2/-4, 1/2, 2/4, 3/6, 3.5/8} {
  \draw (0, \y) -- (-0.08, \y) node[left, font=\scriptsize] {\yt};
}

% Expiratory loop (above axis): rapid rise to PEF, then near-linear
% descent (effort-independent flow limitation)
\draw[curve, engred, smooth] (0.3, 0) .. controls (0.5, 2.5) and (0.8, 3.7) .. (1.1, 3.5)
  .. controls (2.0, 3.2) and (3.5, 1.0) .. (5.0, 0);
\node[engred, font=\small, anchor=west] at (1.2, 3.6) {PEF};
\fill[engred] (1.1, 3.5) circle (1.5pt);

% Inspiratory loop (below axis): symmetric semicircle approximation
\draw[curve, engblue, smooth] (0.3, 0) .. controls (1.5, -2.5) and (3.5, -2.5) .. (5.0, 0);

\node[engred, font=\small, anchor=west] at (3.5, 1.5) {expiration};
\node[engblue, font=\small, anchor=west] at (3.0, -1.5) {inspiration};

\node[font=\scriptsize, anchor=west, align=left] at (4.0, 3.5)
  {FVC, FEV$_1$, PEF};

\node[font=\bfseries, anchor=south] at (2.7, 4.1) {Flow-volume loop};
\end{scope}

\end{tikzpicture}
\caption{Respiratory mechanics. Left: the lung's pressure-volume
relation with the inflation and deflation limbs separated by
hysteresis. The slope at functional residual capacity is the static
compliance, $C \approx 0.2\,\text{L}/\text{cmH}_2\text{O}$ in a
healthy adult. The hysteresis area is the work of breathing per cycle
($\sim 0.4\,\text{J}$ at rest). Right: the spirometric flow-volume
loop with the rapid rise to peak expiratory flow (PEF) at small
expired volume followed by the effort-independent descending limb.
Forced vital capacity (FVC), forced expiratory volume in one second
(FEV$_1$), and PEF are the standard spirometric indices; their
patterns identify obstructive (asthma, COPD; reduced FEV$_1$/FVC) and
restrictive (fibrosis, scoliosis; reduced FVC with preserved
FEV$_1$/FVC) disease.}
\label{fig:vol06ch06:resp-mech}
\end{figure}


\engterm{Quantitative compliance fit}. The static lung P-V curve is
sigmoidal. A Salazar-Knowles two-parameter fit
\[
V(P) = \frac{V_{\max}}{1 + \exp\!\left(-(P - P_{50})/k\right)}
\]
gives, for a healthy adult, $V_{\max} \approx 5.5\,\text{L}$,
$P_{50} \approx 10\,\text{cmH}_2\text{O}$, $k \approx 5\,\text{cmH}_2\text{O}$
\cite[p.~99]{paper:v6c06-salazar-knowles-1964}. The compliance
$C = dV/dP$ is maximum at $P = P_{50}$ where
$C_{\max} = V_{\max}/(4k) \approx 0.28\,\text{L/cmH}_2\text{O}$, and
falls toward $V_{\max}$ at high inflation pressures. The script
\texttt{lung\_compliance\_fit.py} fits the sigmoid and reports the
compliance at FRC ($\approx 0.16\,\text{L/cmH}_2\text{O}$), the
work to inflate a $500\,\text{mL}$ tidal volume above FRC
($\approx 0.36\,\text{J/breath}$), and the work-of-breathing rate at
rest ($\approx 0.07\,\text{W}$) and at exercise
($\approx 0.18\,\text{W}$ at $30$ breaths/min).

\engterm{Surfactant and surface tension}. The lung's alveoli are
lined with a thin liquid film. Surface tension at the air-liquid
interface would, by Laplace's law, collapse the smaller alveoli into
the larger ones unless compensated. The compensation is
\engterm{pulmonary surfactant}, a phospholipid-and-protein mixture
secreted by type II alveolar cells that lowers surface tension and
that decreases tension preferentially as the alveoli compress
\cite[p.~171]{paper:v6c06-clements-1957}. The inverse-area
dependence of surface tension stabilises the alveoli against
collapse. Pre-term infants whose lungs have not yet developed
sufficient surfactant suffer respiratory distress syndrome; the
clinical introduction of surfactant-replacement therapy in the 1990s
dramatically reduced neonatal mortality.

\section{Cell mechanics}\pathtag{standard}

A cell exerts and resists mechanical forces. The mechanical behaviour
of single cells is the subject of cell mechanics, a field that has
matured rapidly with the development of atomic-force microscopy,
optical tweezers, magnetic tweezers, traction-force microscopy, and
micropipette aspiration as routine laboratory tools.

\engterm{Cellular mechanical properties}. A typical animal cell has
a Young's modulus of $100$ to $10\,000\,\text{Pa}$, ranging from
soft (a circulating leukocyte at $\sim 100\,\text{Pa}$) to stiff (a
strongly adherent fibroblast at $\sim 10\,000\,\text{Pa}$). The
stiffness is set by the cortical actin network, the intermediate
filaments, and the cell's internal organelles. Cells under tension
stiffen actively (a phenomenon called active stress stiffening); cells
released from tension soften. The stiffening is mediated by
acto-myosin contraction and provides a feedback loop by which the
cell adapts its stiffness to the mechanical environment.

\engterm{Cell-substrate forces and the force-balance schematic}. A
cell adherent to a substrate exerts traction forces on the substrate
through focal adhesions, the protein complexes that link the cell's
actin cytoskeleton to the substrate via integrin receptors.
Traction-force microscopy measures these forces by tracking the
displacements of beads embedded in an elastic substrate as the cell
pulls. Typical traction stresses for an adherent fibroblast are
$\sim 1\,\text{kPa}$ over a focal adhesion area of
$\sim 1\,\si{\micro\meter}^2$, summing to $\sim 1\,\text{nN}$ per
focal adhesion. A single cell typically has $10$ to $100$ focal
adhesions, summing to $\sim 100\,\text{nN}$ of total traction force.
Figure~\ref{fig:vol06ch06:cell-force} schematises the force balance:
tractions and moments sum to zero at quasi-static equilibrium.

% Figure: Schematic of a single cell on a substrate, with focal
% adhesions, actin stress fibres, and the force balance through the
% integrin-ECM coupling.
\begin{figure}[htbp]
\centering
\begin{tikzpicture}[
  axis/.style={-{Stealth}, thick, draw=engblack},
  ecm/.style={engred!30, draw=engred!70, thick},
  cell/.style={engblue!15, draw=engblue!70, very thick},
  stressfibre/.style={engred, thick},
  fa/.style={engblack, very thick, fill=engblack},
  annot/.style={font=\scriptsize, align=left},
]

% Substrate (ECM)
\draw[ecm, fill=engred!10] (-0.5, -0.5) rectangle (10, 0.5);
\node[annot, anchor=north] at (5.0, -0.6) {Substrate (extracellular matrix, $E_{\text{sub}}$)};

% Cell body: flattened ellipse
\draw[cell] (5, 0.5) ellipse (3.5 and 1.6);

% Nucleus
\draw[engblue!40, draw=engblue!70, thick, fill=engblue!20] (5, 1.0) ellipse (1.0 and 0.5);
\node[font=\scriptsize] at (5, 1.0) {N};

% Focal adhesions (small black rectangles at the basal surface)
\foreach \x in {2.5, 3.4, 4.2, 5.0, 5.8, 6.6, 7.5} {
  \fill[fa] (\x - 0.12, 0.4) rectangle (\x + 0.12, 0.55);
}

% Stress fibres (lines connecting focal adhesions across cell)
\draw[stressfibre] (2.5, 0.5) -- (7.5, 0.5);
\draw[stressfibre] (3.4, 0.5) -- (6.6, 0.5);
\draw[stressfibre] (4.2, 0.5) -- (5.8, 0.5);
\draw[stressfibre] (2.5, 0.5) .. controls (3.5, 1.4) and (6.5, 1.4) .. (7.5, 0.5);

% Traction force arrows: cells pull the substrate inward
\foreach \x/\d in {2.5/-0.5, 3.4/-0.3, 4.2/-0.1, 5.8/0.1, 6.6/0.3, 7.5/0.5} {
  \draw[-{Stealth}, very thick, engblue] (\x, 0.2) -- (\x + \d, 0.2);
}

% Labels
\node[annot, anchor=west] at (8.3, 0.5) {focal\\adhesion};
\draw[-{Stealth}, thin] (8.2, 0.55) -- (7.6, 0.55);

\node[annot, anchor=west] at (8.3, 1.5) {actin stress\\fibre};
\draw[-{Stealth}, thin] (8.2, 1.4) -- (6.5, 0.9);

\node[annot, anchor=east] at (1.5, 1.4) {cell body\\(cortex + cytoplasm)};
\draw[-{Stealth}, thin] (1.6, 1.4) -- (2.5, 1.5);

\node[annot, anchor=west] at (5.0, -0.2) {traction force $\sim 1\,\text{nN}$ per FA};

% Top: stress-fibre tension annotation
\node[annot, engred] at (5.0, 2.4)
  {Stress-fibre tension $T \sim 10$ to $100\,\text{nN}$};

% Sum of traction at focal adhesions
\node[annot, anchor=west] at (-0.5, 3.5)
  {$\displaystyle \sum_i \vec{T}_i = 0$ (cell-substrate force balance)\\
   $\displaystyle \sum_i \vec{T}_i \times \vec{r}_i = 0$ (moment balance)};

\end{tikzpicture}
\caption{Force balance on an adherent cell. The cell exerts traction
forces (blue arrows) on the substrate at focal adhesion plaques,
through transmembrane integrin receptors that link the substrate's
extracellular matrix to the cell's actin cytoskeleton. Internal actin
stress fibres span between focal adhesions and carry tension
$T \sim 10$ to $100\,\text{nN}$ generated by myosin-II motors. The
tractions must sum to zero (force balance) and their moments must
sum to zero (moment balance), because the cell is in quasi-static
equilibrium. Traction-force microscopy measures the tractions by
tracking displacements of beads embedded in an elastic substrate.}
\label{fig:vol06ch06:cell-force}
\end{figure}


\engterm{Mechanotransduction}. Cells convert mechanical signals into
biochemical responses. The principal mechanisms include
mechanically gated ion channels (Piezo1, Piezo2 in mammalian cells),
integrin-mediated signalling at focal adhesions, and stretch-induced
unfolding of cryptic binding sites in cytoskeletal proteins. The
field has shown that cellular phenotype (proliferation,
differentiation, migration, apoptosis) depends on the mechanical
environment in ways that the cell's chemical environment alone does
not explain. The tissue-engineering implication, developed in
Chapter~\ref{vol06:ch05}, is that scaffold stiffness is a control
input on the cell's phenotype.

\section{Soft-tissue constitutive models}\pathtag{mastery}

The constitutive modelling of soft tissues was introduced in
Chapter~\ref{vol06:ch05}. The present section develops the
mathematical detail required to fit specific tissues.

\engterm{The toe-linear-failure schematic}. A collagenous soft tissue
(tendon, ligament, skin, arterial wall) under uniaxial extension
shows three regions: a compliant toe region at small strain where
wavy collagen straightens, a near-linear region at moderate strain
where the straightened fibres load uniformly, and a failure region
where individual fibres begin to fail and the load decreases despite
rising strain. Figure~\ref{fig:vol06ch06:soft-stress-strain} shows
the three regions and their microstructural interpretation.

% Figure: schematic soft-tissue uniaxial stress-strain curve showing
% the three regions (toe, linear, failure) with the underlying
% collagen-microstructure interpretation.
\begin{figure}[htbp]
\centering
\begin{tikzpicture}[
  axis/.style={-{Stealth}, thick, draw=engblack},
  curve/.style={very thick, engred},
  region/.style={draw=none, opacity=0.15},
  annot/.style={font=\scriptsize, align=center},
]

\draw[axis] (0, 0) -- (9, 0) node[right, font=\small] {strain $\varepsilon$};
\draw[axis] (0, 0) -- (0, 6.5) node[above, font=\small] {stress $\sigma$ (MPa)};

\foreach \x/\xt in {1.5/0.02, 3.0/0.04, 4.5/0.06, 6.0/0.08, 7.5/0.10} {
  \draw (\x, 0) -- (\x, -0.08) node[below, font=\scriptsize] {\xt};
}
\foreach \y/\yt in {1/20, 2/40, 3/60, 4/80, 5/100, 5.5/110} {
  \draw (0, \y) -- (-0.08, \y) node[left, font=\scriptsize] {\yt};
}

% Region shading
\fill[engblue, opacity=0.12] (0, 0) rectangle (2.5, 6);
\fill[englime, opacity=0.12] (2.5, 0) rectangle (7.0, 6);
\fill[engred, opacity=0.12] (7.0, 0) rectangle (8.5, 6);

% Toe region: shallow rise
\draw[curve, smooth, samples=40, domain=0:2.5]
  plot (\x, {0.06 * \x * \x});
% Linear region: nearly straight
\draw[curve, smooth, samples=40, domain=2.5:7.0]
  plot (\x, {0.06*2.5*2.5 + 1.15 * (\x - 2.5)});
% Failure region: yield + drop
\draw[curve, smooth, samples=40, domain=7.0:8.0]
  plot (\x, {0.06*2.5*2.5 + 1.15*(7-2.5) + 0.5*(\x - 7) - 5.0*(\x - 7)^2});

% Mark UTS point
\fill[engblack] (7.5, 5.6) circle (2pt);
\node[font=\scriptsize, anchor=south west] at (7.55, 5.65) {UTS};

% Region labels
\node[annot] at (1.25, 5.7) {Toe\\($\varepsilon < 0.02$)};
\node[annot] at (4.75, 5.7) {Linear\\(structural)};
\node[annot] at (7.6, 5.7) {Failure};

% Microstructure cartoons below x-axis
\begin{scope}[yshift=-2.0cm]
% Toe: wavy fibres
\draw[engred, very thick, smooth] (0.3, 0) sin (0.6, 0.15) cos (0.9, 0)
  sin (1.2, -0.15) cos (1.5, 0) sin (1.8, 0.15) cos (2.1, 0);
\node[annot] at (1.25, -0.5) {wavy collagen};
% Linear: straight fibres
\draw[engred, very thick] (3.0, -0.1) -- (6.5, -0.1);
\draw[engred, very thick] (3.0, 0.1) -- (6.5, 0.1);
\node[annot] at (4.75, -0.5) {straightened, stretched};
% Failure: broken
\draw[engred, very thick] (7.0, 0) -- (7.3, 0);
\draw[engred, very thick] (7.5, 0.15) -- (7.7, 0.15);
\draw[engred, very thick] (7.5, -0.15) -- (7.7, -0.15);
\draw[engred, very thick] (7.9, 0) -- (8.2, 0);
\node[annot] at (7.6, -0.5) {fibres failing};
\end{scope}

\end{tikzpicture}
\caption{Schematic uniaxial stress-strain response of a typical
collagenous soft tissue (tendon, ligament, skin). Three regions: the
toe region ($\varepsilon \lesssim 0.02$) where wavy collagen fibres
straighten under small loads with little stiffness; the linear region
($\varepsilon \approx 0.02$ to $0.08$) where the straightened fibres
load uniformly and the modulus is the structural Young's modulus
(typically $0.5$ to $2\,\text{GPa}$ for tendon); and the failure
region where individual fibres begin to fail and the load decreases
even as strain rises. UTS for human Achilles tendon is in the
$50$ to $100\,\text{MPa}$ range. The toe region is the engineering
reason that small tendon deformations cost the body almost no force,
and is the feature that allows tendons to function as efficient
elastic-energy storage elements during locomotion.}
\label{fig:vol06ch06:soft-stress-strain}
\end{figure}


\engterm{Mooney-Rivlin}. For an incompressible isotropic
hyperelastic material with strain-energy density
\[
W = C_{10}(I_1 - 3) + C_{01}(I_2 - 3),
\]
the engineering (first-Piola) stress under uniaxial extension by a
stretch $\lambda$ is
\[
\sigma_{\text{eng}}(\lambda) = 2 \left(\lambda - \lambda^{-2}\right)
   \left(C_{10} + \frac{C_{01}}{\lambda}\right).
\]
Mooney's 1940 derivation and Rivlin's 1948 generalisation laid the
groundwork \cite[p.~588]{paper:v6c06-mooney-1940},
\cite[p.~390]{paper:v6c06-rivlin-1948}. The two-parameter Mooney-Rivlin
model fits rubber-like materials and lightly-strained soft tissues
well, but does not capture the strain-stiffening of collagenous
tissues; for that the Ogden model is the standard.

\engterm{Ogden}. The Ogden 1972 strain-energy function is written in
terms of principal stretches
\cite[p.~568]{paper:v6c06-ogden-1972}:
\[
W = \sum_{p=1}^{N} \frac{\mu_p}{\alpha_p}
\left(\lambda_1^{\alpha_p} + \lambda_2^{\alpha_p} + \lambda_3^{\alpha_p} - 3\right).
\]
For uniaxial extension ($\lambda_1 = \lambda$, $\lambda_2 = \lambda_3 = \lambda^{-1/2}$)
the engineering stress reduces, for $N = 1$, to
\[
\sigma_{\text{eng}}(\lambda) = \frac{\mu}{\alpha}
   \left(\lambda^{\alpha - 1} - \lambda^{-1 - \alpha/2}\right).
\]
The exponent $\alpha$ controls the strain-stiffening: $\alpha = 2$
recovers the neo-Hookean model; $\alpha = 6$ is typical of arterial
wall and skin; $\alpha \gtrsim 10$ captures the stiffer late-region
of tendon. The script \texttt{mooney\_rivlin\_fit.py} fits both
models to a synthetic soft-tissue uniaxial data set and shows that
the Ogden model captures the J-shape stiffening that the Mooney-Rivlin
form misses.

\engterm{Quasi-linear viscoelasticity}. Fung's QLV framework
factorises the soft-tissue stress response into a strain-dependent
elastic part and a strain-independent relaxation function:
\[
\sigma(t) = \int_{-\infty}^{t} G(t - \tau) \frac{\partial \sigma^{(e)}(\varepsilon(\tau))}{\partial \tau} \, d\tau,
\]
where $\sigma^{(e)}(\varepsilon)$ is the elastic-response function
and $G(t)$ is the reduced relaxation function. The framework's
strength is its parsimony: it captures most of the observed
phenomenology with two functions that can be fitted from a handful
of relaxation and step-strain experiments \cite[ch.~7]{text:fung-biomechanics}.

\engterm{Anisotropic models: Holzapfel-Gasser-Ogden}. The
Holzapfel-Gasser-Ogden $2000$ model adds collagen-fibre
contributions to a neo-Hookean matrix
\cite[p.~15]{paper:v6c06-holzapfel-gasser-ogden-2000}:
\[
W = \frac{\mu}{2}(I_1 - 3)
  + \frac{k_1}{2 k_2} \sum_i \left\{
     \exp\!\left[k_2 (\bar{I}_{4i} - 1)^2\right] - 1 \right\},
\]
where $\bar{I}_{4i}$ is the squared stretch in the $i$-th fibre
direction. The HGO model is the working standard for arterial
mechanics and has been used to design aortic stent-grafts,
predict aneurysm rupture, and simulate aortic valve mechanics.

\engterm{Poroelasticity}. Cartilage and other hydrated tissues are
better described as biphasic materials with a solid (collagen and
proteoglycan) phase and a fluid (water) phase that can move relative
to it. The biphasic theory of Mow $1980$
\cite[p.~75]{paper:v6c06-mow-1980} treats cartilage as a porous
elastic solid saturated with an incompressible fluid; load applied
to the cartilage is initially supported by both phases, but fluid
flow out of the loaded region transfers load progressively to the
solid phase over a time scale of seconds to minutes. The biphasic
model explains the time-dependent stress relaxation of cartilage
and is the working framework for cartilage tissue engineering.

\begin{archetype}[Balance, scaling, and failure in biomechanics]
Biomechanics exhibits three of the volume's archetypes in their
canonical applied form. The balance archetype is the force-balance
analysis of body segments by inverse dynamics, the pressure balance
across capillary walls (Starling forces), and the surface-tension
balance in alveoli. The scaling archetype is the allometric
relations across animal sizes (locomotion speed, heart rate,
metabolic rate all scaling with body mass to characteristic powers),
the cubed-radius Murray law for cardiovascular branching, and the
density power-law for trabecular bone modulus. The failure archetype
is the suite of musculoskeletal failure modes (fracture, ligament
rupture, joint degeneration, implant loosening) that ageing and
overuse produce. Biomechanics is the chapter in which the three
archetypes appear together most explicitly.
\end{archetype}

\begin{estimation}
\textbf{Question.} Estimate the peak compressive force on the
patello-femoral joint during a vertical jump from a $40\,\text{cm}$
crouch.

\smallskip
\textit{Estimate, before reading on.}

\smallskip
A $70\,\text{kg}$ adult jumps from a $40\,\text{cm}$ crouch (knee
flexion). The vertical impulse required to lift the body $40\,\text{cm}$
in the air after take-off is
\[
J = m \sqrt{2 g h_{\text{lift}}} = 70 \sqrt{2 \times 9.8 \times 0.4} \approx 196\,\text{kg}\cdot\text{m}/\text{s}.
\]
The impulse must be delivered during the upward push phase, of
duration $\sim 0.3\,\text{s}$. The average upward ground reaction
force during the push is
\[
F_{\text{avg}} = m g + \frac{J}{\Delta t} = 686 + \frac{196}{0.3} \approx 1340\,\text{N}.
\]
Peak ground reaction force is typically $\sim 2\,\text{BW}$ to
$3\,\text{BW}$ (body weights) in vertical jumping, so the peak is
$\sim 2000\,\text{N}$.

The patello-femoral joint reaction force has two contributions: the
ground reaction force passing up through the leg, and the patellar
tendon and quadriceps tendon forces acting at the knee. With the knee
flexed at $\sim 90^{\circ}$ in the deep crouch, the quadriceps
moment arm is $\sim 5\,\text{cm}$ and the external knee-flexion
moment due to ground reaction force (with a moment arm of
$\sim 15\,\text{cm}$ from the ankle) is
$\sim 2000 \times 0.15 = 300\,\text{N}\cdot\text{m}$. The quadriceps
must produce $300/0.05 = 6000\,\text{N}$ of tendon tension. The
patello-femoral compression force is approximately $1.5$ times the
patellar tendon tension at this knee angle, or $\sim 9000\,\text{N}$
\cite[p.~1078]{paper:v6c06-pellegrino-2003}.

\smallskip
\textit{Postmortem.} Measured peak patello-femoral joint forces
during a vertical jump (from inverse-dynamics analyses) are
$\sim 8000$ to $12\,000\,\text{N}$ for adult subjects, or $11$ to
$17\,\text{BW}$. The estimate is at the low end of the measured
range, which is appropriate given the simplifications (we did not
include the quadriceps inertia contribution and used a small-angle
estimate for the moment-arm geometry). The remarkable observation is
how large the joint force is: nine kilonewtons on a contact area of
$\sim 4\,\text{cm}^2$ produces an average joint pressure of
$\sim 22\,\text{MPa}$. The compression is at the boundary of what
articular cartilage can repeatedly tolerate without damage, which is
the engineering reason that explosive jumping is a common source of
patello-femoral pain and cartilage degeneration in athletes.
\end{estimation}

\begin{estimation}
\textbf{Question.} Estimate the cardiac mechanical power output of a
$70\,\text{kg}$ adult cycling at $250\,\text{W}$ mechanical output,
and compare to the resting cardiac power.

\smallskip
\textit{Estimate, before reading on.}

\smallskip
At rest, cardiac output is $\dot Q \approx 5\,\text{L/min} =
8.3 \times 10^{-5}\,\text{m}^3/\text{s}$ and mean arterial pressure
is $\bar P \approx 100\,\text{mmHg} \approx 1.33\times 10^{4}\,\text{Pa}$.
The left-ventricular power is
\[
P_{\text{LV,rest}} \approx \bar P\, \dot Q \approx 1.33\times 10^{4}
\times 8.3\times 10^{-5} \approx 1.1\,\text{W}.
\]
Add the kinetic-energy term $\tfrac{1}{2}\rho v^2 \dot Q$ for the
$\sim 17\,\text{cm/s}$ aortic mean velocity ($\sim 0.02\,\text{W}$,
negligible) and the right-ventricular power
($\sim 0.2\,\text{W}$ at $20\,\text{mmHg}$ pulmonary mean), for a
total cardiac power $\sim 1.3\,\text{W}$ at rest.

At $250\,\text{W}$ mechanical output, a trained adult has
$\dot V_{\text{O}_2} \approx 3.0\,\text{L/min}$, so cardiac output
rises to $\dot Q \approx 20\,\text{L/min} = 3.3\times 10^{-4}\,\text{m}^3/\text{s}$
and mean arterial pressure to $\bar P \approx 130\,\text{mmHg}$. The
cardiac power is
\[
P_{\text{LV,ex}} \approx 1.73\times 10^{4}\,\text{Pa} \cdot
3.3\times 10^{-4}\,\text{m}^3/\text{s} \approx 5.7\,\text{W}.
\]
Add the pulmonary side and the velocity term:
$P_{\text{cardiac,ex}} \approx 7\,\text{W}$.

\smallskip
\textit{Postmortem.} Measured cardiac mechanical power during heavy
cycling is $5$ to $10\,\text{W}$, in line with the estimate. The
heart is metabolically inefficient ($\sim 25\,\%$): a $7\,\text{W}$
output corresponds to a $28\,\text{W}$ myocardial oxygen demand,
which is the engineering reason that heavy exercise is rate-limited
by coronary blood flow in the presence of obstructive coronary
disease.
\end{estimation}

\section{Failure: musculoskeletal injury and orthopaedic implants}\pathtag{core}

The musculoskeletal system fails in well-characterised mechanism
classes that occupy the bulk of orthopaedic clinical practice.

\engterm{Bone fracture}. A bone fractures when the applied stress
exceeds the bone's strength. Fracture types include
\engterm{transverse} (perpendicular to the long axis, from bending
or direct impact), \engterm{spiral} (oblique helical, from torsion),
\engterm{comminuted} (multiple fragments, from high-energy impact),
and \engterm{stress fracture} (small crack from repeated submaximal
loading). The principal age-associated bone-fracture mechanism is
osteoporotic fragility: the loss of bone mass with age
($\sim 1\,\%/\text{year}$ after age 30, accelerating in women after
menopause) reduces the bone's load-bearing capacity to the point
where ordinary activities (a low-energy fall, a cough, lifting
groceries) can cause fractures. Osteoporotic fractures of the hip,
wrist, and vertebra affect approximately $30\,\%$ of women and
$15\,\%$ of men over age 65 (current as of 2024) and constitute a
major contributor to disability and mortality in the elderly.

\engterm{Stress fracture mechanics}. Stress fractures arise from
fatigue, the same accumulation-of-microcracks mechanism that drives
metal fatigue (Volume X). The bone's cyclic strength is well below
its monotonic strength: at $\sim 4000$ cycles, cortical bone fails at
$\sim 100\,\text{MPa}$ longitudinal compressive amplitude, half its
monotonic value; at $\sim 10^7$ cycles, the endurance limit is
$\sim 30\,\text{MPa}$. The running impact of a $70\,\text{kg}$
adult delivers $\sim 30\,\text{MPa}$ peak stress in the tibial
mid-shaft cortex, which is the engineering reason that high-mileage
runners regularly develop tibial stress fractures: each running step
is a low-amplitude fatigue cycle, and $10\,000$ steps per day
accumulate $\sim 4 \times 10^6$ cycles per year before any
remodelling repair.

\engterm{Ligament and tendon injury}. Ligaments and tendons fail
under excessive tension. The anterior cruciate ligament (ACL) of the
knee is the most commonly torn ligament in sports medicine; the
mechanism is a non-contact rotational injury at the knee with the
foot planted, producing peak ligament forces above the ACL's tensile
strength ($\sim 2000\,\text{N}$). Boden et al.\ 2000
\cite[p.~575]{paper:v6c06-tabin-acl-1996} identify the canonical
mechanism: anterior tibial translation combined with internal
rotation and valgus moment at the knee, often during a deceleration
or cutting manoeuvre. Tendon injuries include tendinopathy (chronic
overuse with collagen degeneration but no acute rupture), partial
tear, and complete rupture (most commonly of the Achilles,
supraspinatus, and patellar tendons). The biomechanics of these
injuries are now well-understood; the treatment ranges from
conservative rehabilitation to surgical repair or reconstruction.

\engterm{Rotator cuff and shoulder injury}. The supraspinatus tendon,
which inserts on the greater tuberosity of the humerus through the
subacromial space, is one of the most commonly torn tendons in
adults over age 50. The mechanism is a combination of degenerative
tendon weakening (collagen fibre disorganisation, reduced
cellularity) and mechanical impingement against the coracoacromial
arch during arm elevation. Surgical repair restores the anatomy in
most cases but the failure mode persists in the contralateral
shoulder, suggesting that the dominant cause is systemic rather
than local mechanical.

\engterm{Joint degeneration: osteoarthritis}. Osteoarthritis is the
progressive degradation of articular cartilage, leading to joint
pain, stiffness, and loss of function. The mechanism combines
mechanical wear (loss of proteoglycan content, fibrillation of the
cartilage surface), biological catabolism (matrix metalloproteinases
secreted by the chondrocytes degrade the matrix), and reduced
synthesis (chondrocyte senescence reduces matrix renewal)
\cite[p.~1241]{paper:v6c06-pelletier-2001}. Risk factors include age,
obesity, joint injury, and joint malalignment; the disease affects
approximately $20\,\%$ of adults over age 60 (current as of 2024) and
is the leading cause of joint replacement.

\engterm{Orthopaedic-implant failure}. Joint replacements (hip, knee,
shoulder, ankle) and spinal hardware can fail by several
characteristic mechanisms. \engterm{Aseptic loosening} is the
progressive separation of the implant from the surrounding bone due
to bone resorption at the bone-implant interface; the resorption is
triggered by wear-particle-induced inflammation and is the dominant
late failure mode of hip and knee replacements. \engterm{Periprosthetic
infection} is bacterial colonisation of the implant surface, often
producing a biofilm that resists antibiotic treatment and that
requires implant removal. \engterm{Mechanical failure} of the
implant itself (fatigue fracture of femoral stems, polyethylene
liner wear, ceramic head fracture) is now rare with current
materials but was historically a major failure mode. \engterm{Bearing
surface wear} produces submicrometre wear particles that the immune
system recognises as foreign and that drive the aseptic-loosening
cascade.

\engterm{Hip-replacement long-term outcomes}. Charnley's $1972$ paper
on the low-friction arthroplasty established the modern hip
replacement: a small ($22\,\text{mm}$) metal femoral head, an
ultra-high-molecular-weight polyethylene (UHMWPE) acetabular cup,
and PMMA bone cement to fix the prosthesis in the bone
\cite[p.~64]{paper:v6c06-charnley-1972}. The current revision rate
at $20$ years after primary hip replacement is $\sim 10\,\%$
(current as of 2024), down from $\sim 30\,\%$ in the $1980$s
implants \cite[ch.~14]{text:mow-huiskes-orthopaedic}. The
hip-replacement projection literature
\cite[Tab.~1]{paper:v6c06-kurtz-2007} estimates that primary hip
replacements in the US will rise from $\sim 200\,000$/year in $2005$
to $\sim 570\,000$/year by $2030$ as the population ages. Pauwels
$1976$ analysis of the hip-joint force balance (single-leg stance:
the gluteus medius must produce $\sim 2.5\,\text{BW}$ of tendon
tension to balance the body weight at the offset hip centre)
remains the working model for prosthesis design
\cite[ch.~3]{paper:v6c06-pauwels-1976}.

The engineering response to each failure mode is well-characterised:
modern hip and knee replacements use cross-linked polyethylene
bearings with much lower wear rates, ceramic-on-ceramic alternatives
for high-demand patients, antibiotic-impregnated bone cement for
infection prevention, and porous-coated implant surfaces to encourage
biological fixation through bone in-growth.

\begin{principle}[Musculoskeletal failure has predictable mechanism classes]
The musculoskeletal system fails by a small set of well-characterised
mechanism classes: bone fracture (acute, fatigue, fragility),
ligament-tendon rupture (acute, chronic), joint degeneration
(osteoarthritis), and implant failure (loosening, infection,
wear). Each mechanism class has well-characterised risk factors,
predictable progression, and corresponding engineering responses
(implant design, surgical technique, rehabilitation protocol). The
clinician or biomechanical engineer assesses a presenting injury or
device failure by first identifying the mechanism class and then
applying the corresponding response.
\end{principle}

\section{Project}\pathtag{core}

\begin{project}[Instrument a vertical jump and estimate peak knee force]
The reader instruments a series of vertical jumps with a smartphone
accelerometer or inexpensive force sensor and computes the peak knee
force by inverse dynamics.

\textbf{Track.} Physical experiment plus analysis.
\textbf{Hazard class}: low. The subject should be in normal health
and a regular participant in jumping activity. Avoid the experiment
if there is recent lower-limb injury.

\textbf{Equipment.} A smartphone with accelerometer logging app, or
an inexpensive USB force plate or load cell; a tape measure; a
metronome or stopwatch.

\textbf{Procedure.}

\begin{enumerate}[leftmargin=*]
\item Attach the smartphone securely to the subject's lower back
  (the approximate centre of mass) with a belt or strap. Calibrate
  the accelerometer's orientation.
\item Have the subject perform $10$ vertical jumps from a
  $\sim 30\,\text{cm}$ crouch, with $30\,\text{seconds}$ rest between
  jumps. Log the accelerometer data continuously.
\item For each jump, identify the take-off and landing peaks in the
  acceleration data. Compute the impulse during the take-off push
  phase by integrating the upward acceleration over the push duration.
\item Compute the average and peak ground reaction force during the
  push phase from the impulse and the subject's mass.
\item Compute the peak patello-femoral joint force using the
  geometric model in the estimation block (or a more detailed
  inverse-dynamics model from a standard biomechanics text).
\item Repeat with three different crouch depths (shallow $\sim 15\,\text{cm}$,
  medium $\sim 30\,\text{cm}$, deep $\sim 50\,\text{cm}$). Compare
  the peak forces and the jump heights.
\end{enumerate}

\textbf{Deliverable.} A ten- to fifteen-page report with one
acceleration-time figure (representative jump), one table summarising
the ten jumps' peak forces and jump heights, one plot of peak knee
force versus crouch depth, and a closing paragraph naming the largest
source of uncertainty in the estimate (sensor placement, soft-tissue
artifact, segment-mass estimates, geometric model).

\textbf{Time.} Two to four hours of measurement; three to six hours
of analysis.

\textbf{Reference code.} \texttt{jump\_peak\_force.py} in this
chapter's \texttt{code/} directory implements the impulse computation
and the knee-force geometric model; the reader who does not wish to
collect data can run the synthetic-jump example.
\end{project}

\section{Exercises}

\subsection*{Calculation}\pathtag{core}

\begin{exercise}
Cortical bone has Young's modulus of $15\,\text{GPa}$ along the long
axis and compressive strength of $200\,\text{MPa}$. Compute the
elastic strain at failure under uniaxial compression.
\end{exercise}

\paragraph{Solution sketch.}
For a linear-elastic response up to failure,
$\varepsilon_{\text{f}} = \sigma_{\text{f}}/E = 200\,\text{MPa}/15\,\text{GPa}
\approx 1.33 \times 10^{-2}$, i.e.\ approximately $1.3\,\%$ strain at
the compressive yield. Measured failure strain in human cortical
bone is closer to $2\,\%$, the discrepancy reflecting a small
post-yield plastic phase before final fracture
\cite[Tab.~3]{paper:v6c06-reilly-burstein-1975}. The strain is
intermediate between concrete ($\sim 0.3\,\%$) and ductile steel
($\sim 20\,\%$): cortical bone is a brittle composite at the
structural scale, with five to ten times more strain-to-failure
than unreinforced ceramic.

\begin{exercise}
A human Achilles tendon has cross-sectional area $\sim 80\,\text{mm}^2$,
length $\sim 15\,\text{cm}$, and Young's modulus $\sim 1\,\text{GPa}$.
Compute the elongation under a tensile load of $4000\,\text{N}$
(typical peak load during running) and the elastic energy stored.
\end{exercise}

\paragraph{Solution sketch.}
Stress $\sigma = F/A = 4000/(80 \times 10^{-6}) = 50\,\text{MPa}$.
Strain $\varepsilon = \sigma/E = 50\,\text{MPa}/1\,\text{GPa} = 0.05
= 5\,\%$. Elongation $\Delta L = \varepsilon L = 0.05 \times
0.15\,\text{m} = 7.5\,\text{mm}$. Stored elastic energy
$U = \tfrac{1}{2} F \Delta L = \tfrac{1}{2} \times 4000 \times 0.0075
= 15\,\text{J}$. For two Achilles tendons working alternately at
roughly $30\,\%$ return efficiency, the per-stride energy returned is
$\sim 30 \times 2 \times 0.3 \approx 18\,\text{J}$, comparable to the
$\sim 50\,\text{J}/\text{stride}$ total mechanical work of running,
which is the quantitative explanation for the $\sim 35\,\%$ metabolic
saving of tendon elasticity in running
\cite[p.~64]{paper:v6c06-alexander-1991}.

\begin{exercise}
The heart pumps $5\,\text{L}/\text{min}$ against a mean pressure of
$100\,\text{mmHg}$. Compute the cardiac power output in watts.
\end{exercise}

\paragraph{Solution sketch.}
$\dot Q = 5\,\text{L/min} = 8.33 \times 10^{-5}\,\text{m}^3/\text{s}$;
$\bar P = 100\,\text{mmHg} = 100 \times 133.322 \approx 13332\,\text{Pa}$.
Power $P = \bar P \dot Q = 13332 \times 8.33 \times 10^{-5}
\approx 1.11\,\text{W}$. Add the pulmonary side
($\sim 20\,\text{mmHg} \times 8.33\times 10^{-5}\,\text{m}^3/\text{s}
\approx 0.22\,\text{W}$) for a total left + right ventricular
mechanical power $\sim 1.3\,\text{W}$ at rest, which is $\sim 7\,\%$
of the $\sim 20\,\text{W}$ resting metabolic rate of a typical adult
heart (myocardial efficiency $\sim 25\,\%$).

\begin{exercise}
The pulse-wave velocity in a healthy artery is $\sim 7\,\text{m}/\text{s}$.
Use the Moens-Korteweg equation to estimate the arterial wall's
incremental elastic modulus, given $h = 1\,\text{mm}$,
$\rho_{\text{blood}} = 1060\,\text{kg}/\text{m}^3$, $R = 10\,\text{mm}$.
\end{exercise}

\paragraph{Solution sketch.}
$c = \sqrt{E h / (2 \rho R)}$, so
$E = 2 \rho R c^2 / h
   = 2 \times 1060 \times 10^{-2} \times 49 / 10^{-3}
   \approx 1.04 \times 10^{6}\,\text{Pa}
   = 1.0\,\text{MPa}$.
The arterial wall's incremental modulus is in the low-MPa range, in
line with measured aortic incremental moduli of $0.4$ to $1.5\,\text{MPa}$
in healthy adults. At PWV $= 10\,\text{m/s}$ (seventy-year-old) the
inferred modulus rises to $\sim 2.1\,\text{MPa}$, doubling the wall
stiffness, which is the engineering basis for using PWV as a
clinical biomarker of arterial ageing.

\begin{exercise}
A typical lung has compliance $C = 200\,\text{mL}/\text{cmH}_2\text{O}$
and airway resistance $R = 1\,\text{cmH}_2\text{O}/(\text{L}/\text{s})$.
Compute the time constant $RC$ for passive expiration and compare to
the typical breath duration of $\sim 4\,\text{s}$.
\end{exercise}

\paragraph{Solution sketch.}
Unit-consistent product: $RC = 1\,\text{cmH}_2\text{O}/(\text{L/s})
\times 0.2\,\text{L/cmH}_2\text{O} = 0.2\,\text{s}$. The time constant
is $\sim 5\,\%$ of the breath duration, so passive expiration is
essentially complete within one third of the breath cycle, leaving
the rest of the expiratory phase as a brief pause before the next
inspiration. In COPD the resistance rises (to $5\text{-}10\,\text{cmH}_2\text{O}/(\text{L/s})$),
pushing $RC$ to $1\text{-}2\,\text{s}$, the same order as the breath
period; the consequence is dynamic hyperinflation (incomplete
expiration), the clinical signature of obstructive disease.

\begin{exercise}
A capillary has radius $4\,\si{\micro\meter}$, length $1\,\text{mm}$,
and blood viscosity $4\,\text{cP}$. Compute the pressure drop required
to drive blood through the capillary at a typical capillary flow
velocity of $0.05\,\text{cm}/\text{s}$.
\end{exercise}

\paragraph{Solution sketch.}
Use Poiseuille:
$\Delta P = 8 \mu L \bar v / R^2$
with $\mu = 4 \times 10^{-3}\,\text{Pa}\cdot\text{s}$,
$L = 10^{-3}\,\text{m}$, $\bar v = 5\times 10^{-4}\,\text{m/s}$,
$R = 4\times 10^{-6}\,\text{m}$:
\[
\Delta P = \frac{8 \cdot 4\times 10^{-3} \cdot 10^{-3} \cdot 5\times 10^{-4}}
                {(4\times 10^{-6})^2}
        = \frac{1.6\times 10^{-8}}{1.6\times 10^{-11}}
        = 1000\,\text{Pa} \approx 7.5\,\text{mmHg}.
\]
The arteriole-to-venule pressure drop across the capillary bed is
typically $\sim 15\,\text{mmHg}$ (arteriole exit $\sim 30$,
venule entry $\sim 15$); the calculation accounts for half of that
drop in a single capillary path, in line with the typical
two-capillary-in-series anatomy. The Fahraeus-Lindqvist effect
reduces the effective viscosity at this radius by $\sim 30\,\%$,
which would lower the calculated $\Delta P$ to $\sim 5\,\text{mmHg}$.

\begin{exercise}
A cell exerts a traction force of $1\,\text{nN}$ at each of $50$
focal adhesions. Compute the total cellular traction force and
compare to the cell's weight (mass $\sim 10^{-12}\,\text{kg}$).
\end{exercise}

\paragraph{Solution sketch.}
The tractions are oriented inward at each focal adhesion; if the
cell is in quasi-static equilibrium the vector sum is zero, but the
sum of magnitudes is $50 \times 1\,\text{nN} = 50\,\text{nN}$. The
cell's weight is $m g \approx 10^{-12} \times 9.8 \approx
10^{-11}\,\text{N} = 10\,\text{pN}$. The traction-force magnitude is
$\sim 5000$ times the cell's weight: cells live in the inertia-free,
gravity-irrelevant regime, and the dominant mechanical forces are
generated by the cell's own actomyosin contraction at focal
adhesions, not by external loads.

\subsection*{Derivation}\pathtag{standard}

\begin{exercise}
Derive the Moens-Korteweg equation for pulse-wave velocity in a
thin-walled elastic tube containing an incompressible fluid. State
the assumptions and identify when each fails.
\end{exercise}

\paragraph{Solution sketch.}
Take a thin-walled elastic tube of radius $R$, wall thickness $h$,
elastic modulus $E$, containing fluid of density $\rho$. A small
pressure perturbation $\delta P$ produces a wall distension
$\delta R$ found from the hoop-stress balance
$\delta P \cdot R = E h \cdot \delta R / R$ (thin-wall
approximation: hoop stress $\sigma_{\theta} = P R / h$, strain
$\delta R / R$). Hence
$\delta R = \delta P \cdot R^2 / (E h)$.
The cross-section area is $A = \pi R^2$, so
$\delta A / A = 2 \delta R / R = 2 \delta P R / (E h)$.

Linearised mass conservation in a $1$-D fluid:
$\partial A/\partial t + \partial (A u)/\partial x = 0$; linearised
momentum: $\rho \partial u/\partial t = -\partial P/\partial x$.
Combining and using $\delta A = A_0 \cdot 2\delta P R/(Eh)$ gives a
wave equation for $P$ with wave speed
\[
c^2 = \frac{E h}{2 \rho R}.
\]
Assumptions and failure modes: \emph{thin wall} ($h \ll R$; fails at
the great arteries where $h/R \sim 0.1$, requiring the Korteweg
correction); \emph{linear elasticity} (fails at large pulse pressures
where the artery enters the J-region; the apparent PWV rises with
mean pressure); \emph{incompressible fluid} (blood is essentially
incompressible); \emph{negligible viscous losses} (fails for small
arteries where Womersley number $< 1$ and viscous boundary layer
fills the tube). The clinical interpretation: in a stiffening
artery the modulus $E$ rises and PWV rises with it; PWV measurement
is the non-invasive surrogate.

\begin{exercise}
Derive the Hill muscle force-velocity relation
$(F + a)(v + b) = (F_0 + a) b$ from the assumption that the
energy released by the muscle exceeds the mechanical work by an
amount that increases linearly with shortening velocity. Identify
the physiological meaning of the parameters.
\end{exercise}

\paragraph{Solution sketch.}
Hill observed that the rate of energy liberation in a shortening
muscle is the sum of the mechanical work rate ($Fv$) and an
``extra heat'' proportional to the shortening velocity ($av$):
$dE/dt = Fv + av$. He also observed that the total energy rate
declines with force: as force rises, the velocity falls, and the
energy rate is observed to follow $dE/dt = b(F_0 - F)$ at constant
$F_0$. Setting the two expressions equal,
\[
Fv + av = b(F_0 - F) \quad\Longrightarrow\quad
(F + a)(v + b) = (F_0 + a) b
\]
after rearrangement \cite[p.~151]{paper:v6c06-hill-1938}.
Physiological meaning: $F_0$ is the isometric peak force (at $v = 0$);
$v_{\max} = b F_0/a$ is the maximum shortening velocity (force at
zero load); $a$ has units of force and is the rate of heat
liberation per unit shortening; $b$ has units of velocity and is
the rate constant of the contraction. The empirical values
$a/F_0 \approx 0.25$ and $b/v_{\max} \approx 0.25$ produce the
characteristic hyperbola illustrated in
Figure~\ref{fig:vol06ch06:muscle-curves}.

\begin{exercise}
Derive the Laplace law for alveolar surface tension
$P = 2\gamma/r$ for a spherical alveolus, and explain why the
inverse dependence on radius would, without surfactant, destabilise
the smaller alveoli into the larger ones.
\end{exercise}

\paragraph{Solution sketch.}
For a thin spherical liquid film (single liquid-air interface, as in
an alveolus, not a soap bubble's double interface) of radius $r$ and
surface tension $\gamma$, the force balance is between the
pressure-induced outward force and the surface-tension hoop force.
Consider the upper hemisphere of the alveolus: the outward
pressure-area force is $\Delta P \cdot \pi r^2$; the surface-tension
force along the equatorial circle is $\gamma \cdot 2\pi r$. Setting
them equal,
$\Delta P = 2\gamma/r$.

For a closed alveolar tree where two alveoli share a common airway
duct, the smaller alveolus has the higher internal pressure
($\propto 1/r$). The pressure gradient drives air from small to
large, collapsing the small alveolus and over-distending the large
one, the well-known Laplace instability. Pulmonary surfactant
breaks the instability by making $\gamma$ depend on the surface
density of the phospholipid film: at small $r$ (small alveolus,
high surface density), $\gamma$ falls preferentially, producing a
near-constant $\Delta P$ across alveoli of different size
\cite[p.~171]{paper:v6c06-clements-1957}.

\subsection*{Estimation}\pathtag{core}

\begin{exercise}
Estimate the metabolic cost of walking $1\,\text{km}$ at the optimum
walking speed for a $70\,\text{kg}$ adult. Compare to the cost of
running the same distance.
\end{exercise}

\paragraph{Solution sketch.}
At optimum walking, metabolic cost is $\sim 0.5\,\text{J}/(\text{kg}\cdot\text{m})$.
For a $70\,\text{kg}$ adult over $1\,\text{km}$:
$E_{\text{walk}} \approx 0.5 \times 70 \times 1000 = 35\,\text{kJ}
\approx 8.4\,\text{kcal}$, or $\sim 50\,\text{kcal}$ for a typical
$5\,\text{km}$ walk. Running cost is approximately
$1.0\,\text{J}/(\text{kg}\cdot\text{m})$, roughly twice the walking
cost, so $E_{\text{run}} \approx 70\,\text{kJ} \approx
17\,\text{kcal}/\text{km}$. The running-vs-walking energy ratio is
approximately constant with distance because both costs are
dominated by per-stride mechanical work and not by aerobic baseline.
The often-quoted heuristic ``$100\,\text{kcal/mile}$'' is
approximately right for running at $1$ mile $\approx 1.6\,\text{km}$.

\begin{exercise}
Estimate the peak force in the Achilles tendon during a heel-toe
landing from a $50\,\text{cm}$ jump. Compare to the tendon's
ultimate tensile strength.
\end{exercise}

\paragraph{Solution sketch.}
Landing from $50\,\text{cm}$ delivers vertical velocity
$v = \sqrt{2 g h} = 3.1\,\text{m/s}$ at touchdown. Take a typical
landing deceleration time of $0.1\,\text{s}$ for a heel-toe
landing (longer than for a stiff landing). Mean deceleration:
$\bar a = v/\Delta t = 31\,\text{m/s}^2 \approx 3\,g$; peak
$\sim 2\bar a \approx 6\,g$. Mean vertical GRF: $\sim 4\,\text{BW}
\approx 2800\,\text{N}$; peak $\sim 6\text{-}8\,\text{BW} \approx
5000\,\text{N}$. During the landing the ankle plantar-flexors
resist dorsiflexion; the moment arm of the Achilles tendon at the
ankle is $\sim 5\,\text{cm}$ and the moment arm of the GRF is
$\sim 10\,\text{cm}$ (forefoot loading); the Achilles tendon
tension is approximately $2 \times$ GRF $\approx 10\,000\,\text{N}$.

The Achilles' UTS is $\sim 100\,\text{MPa}$ at $\sim 80\,\text{mm}^2$
cross-section, giving an ultimate load of $\sim 8000\,\text{N}$. The
estimate puts the landing load at or above ultimate, which is the
quantitative reason that Achilles rupture is a frequent injury in
stiff landings from height; the actual safety margin is recovered by
the tendon's strain-rate stiffening and the muscle's pre-activation
that distributes the load over a longer time window.

\begin{exercise}
Estimate the time required for an oxygen molecule to diffuse from
the alveolar lumen into the pulmonary capillary blood across the
$\sim 1\,\si{\micro\meter}$ alveolar wall.
\end{exercise}

\paragraph{Solution sketch.}
Use $t \sim L^2/D$. For oxygen in water (and the alveolar membrane
is mostly water with embedded protein), $D \approx 2\times 10^{-9}\,\text{m}^2/\text{s}$.
With $L = 10^{-6}\,\text{m}$:
\[
t \approx \frac{(10^{-6})^2}{2 \times 10^{-9}} = 5 \times 10^{-4}\,\text{s} = 0.5\,\text{ms}.
\]
The capillary transit time is $\sim 0.75\,\text{s}$ at rest and falls
to $\sim 0.25\,\text{s}$ at heavy exercise. The diffusion time is
three orders of magnitude faster, so oxygenation is normally
diffusion-rich (not diffusion-limited) at rest and only becomes
limited at extreme exercise when transit time approaches the
diffusion time. In pulmonary fibrosis, where the membrane thickness
doubles, diffusion time rises quadratically ($\times 4$) and
becomes the bottleneck during exertion.

\begin{exercise}
Estimate the volume of blood in the body's capillaries, given a
typical capillary cross-section of $50\,\si{\micro\meter}^2$ and a
total capillary length of $\sim 60\,000\,\text{km}$.
\end{exercise}

\paragraph{Solution sketch.}
$V = A \times L = 50 \times 10^{-12}\,\text{m}^2 \times 6 \times 10^{7}\,\text{m}
= 3 \times 10^{-3}\,\text{m}^3 = 3\,\text{L}$. This is comparable to
the body's total blood volume of $\sim 5\,\text{L}$, but well above
the typical statement that ``capillaries hold $\sim 5\,\%$ of blood
at any moment''. The discrepancy is because not all capillaries are
perfused simultaneously: typical resting capillary recruitment is
$\sim 30\text{-}50\,\%$ of the anatomic bed, with the remainder
shunted through arteriovenous anastomoses or simply closed. The
estimate gives the anatomical reserve; the functional volume is
$\sim 0.3\,\text{L}$ at rest, rising during exercise.

\subsection*{Simulation}\pathtag{mastery}

\begin{exercise}
Implement a finite-element model of an adult tibia under axial
compression. Compute the strain distribution and identify the
location of peak strain.
\end{exercise}

\paragraph{Discussion.}
Rubric: (i) Geometry. Model the tibia as a tapered hollow cylinder
with $30\,\text{mm}$ proximal radius, $15\,\text{mm}$ distal radius,
$5\,\text{mm}$ cortical wall thickness, $40\,\text{cm}$ length.
Mesh with $\sim 10^4$ tetrahedral or hexahedral elements.
(ii) Material model. Use cortical bone $E = 18\,\text{GPa}$,
$\nu = 0.4$; trabecular bone $E = 1\,\text{GPa}$, $\nu = 0.3$;
either ignore the marrow cavity or assign $E = 1\,\text{MPa}$.
(iii) Loading. Apply $2500\,\text{N}$ (running peak axial load,
$\sim 3.5\,\text{BW}$) on the proximal articular surface;
fix the distal articular surface.
(iv) Results. Expected peak strain at the medial cortex of the
mid-shaft, $\sim 2500\,$microstrain for the running load, within
the safe-stress envelope. Peak strain at the postero-medial
mid-shaft is where tibial stress fractures concentrate in runners.
(v) Validation. Compare to published in-vivo strain-gauge
measurements (Lanyon \& Hartman 1979; Burr 1996), which report
$1500\text{-}3000\,$microstrain on the medial tibia during running.

\begin{exercise}
Simulate a Windkessel model of the arterial system as a two-element
RC circuit. Compute the systolic-to-diastolic pressure decay during
diastole and compare to typical measured pressure waveforms.
\end{exercise}

\paragraph{Solution sketch.}
The two-element Windkessel equation is
$C\,dP/dt + P/R = Q(t)$, where $C$ is total arterial compliance
($\sim 1.5\,\text{mL/mmHg}$ in an adult) and $R$ is total peripheral
resistance ($\sim 17\,\text{mmHg/(L/min)}$). During diastole $Q = 0$,
so $P(t) = P_{\text{end-systole}} \cdot \exp(-t/RC)$, with the
time constant $RC$. Unit-consistent conversion gives
$RC \approx 1.5\,\text{s}$ for the standard parameter set.
Diastolic duration is $\sim 0.6\,\text{s}$ at $70\,\text{bpm}$, so
$P_{\text{end-diastole}}/P_{\text{end-systole}}
= \exp(-0.6/1.5) \approx 0.67$, i.e.\ pressure falls from
$\sim 120$ to $\sim 80\,\text{mmHg}$, matching the typical $120/80$
measurement. The model fails for the dicrotic notch (a
high-frequency reflection phenomenon) and for the wave-reflection
-driven systolic peak; a four-element extension (adding aortic input
impedance and a proximal compliance) captures both.

\subsection*{Design}\pathtag{standard}

\begin{exercise}
Design a femoral component for a total hip replacement. Specify the
material, geometry, and surface finish. Defend the choice against
the alternatives (stainless steel versus titanium versus cobalt-chrome;
cemented versus press-fit; smooth versus porous-coated).
\end{exercise}

\paragraph{Discussion.}
Rubric: (i) Material. Titanium alloy (Ti-6Al-4V or Ti-6Al-7Nb) for
the stem because of its lower modulus ($110\,\text{GPa}$ vs cortical
bone's $18\,\text{GPa}$; cobalt-chrome at $230\,\text{GPa}$ produces
worse stress shielding), excellent fatigue strength, and corrosion
resistance. Cobalt-chrome head (CoCrMo) for wear resistance against
the polyethylene cup. Stainless steel rejected because of fatigue
and corrosion concerns in long-term implantation.
(ii) Geometry. Anatomically contoured stem (rather than straight)
that fits the proximal femoral canal; $35$-$45\,\text{mm}$
collar-to-tip length; $10$-$15\,\text{mm}$ neck length offset matched
to the patient's femoral offset. The Pauwels analysis
\cite[ch.~3]{paper:v6c06-pauwels-1976} sets the load axis: the
component must carry $\sim 2.5\,\text{BW}$ at the offset hip centre.
(iii) Fixation. Porous-coated press-fit for patients under 65 (long
expected life, biological fixation through bone ingrowth); cemented
for patients over 75 with poor bone quality (better short-term
stability, avoiding peri-prosthetic fracture). Hybrid (cemented stem,
press-fit cup) is the current US standard.
(iv) Surface. Porous coating (CoCrMo beads or Ti plasma spray) on the
proximal third of the stem; smooth polished distal stem to reduce
debris generation if there is micromotion.
(v) Trade-off discussion. Cemented gives early stability but lifetime
limited by cement-bone interface; porous-coated requires careful
press-fit but achieves biological fixation that lasts. Charnley's
$1972$ original design was cemented metal-on-UHMWPE
\cite[p.~64]{paper:v6c06-charnley-1972}; the modern revision rate
is one-third of the $1980$s rate
\cite[ch.~14]{text:mow-huiskes-orthopaedic} due to cross-linked
polyethylene, ceramic options, and improved porous-coating
techniques.

\begin{exercise}
Design a prosthetic running blade for a below-knee amputee. Specify
the material (carbon-fibre composite), the geometric shape (J-shaped
curve), and the energy-return characteristics required.
\end{exercise}

\paragraph{Discussion.}
Rubric: (i) Material. Unidirectional carbon-fibre / epoxy composite
($\sim 60\,\%$ fibre volume), $E_{\text{long}} \approx 130\,\text{GPa}$,
specific strength higher than steel and titanium, fatigue strength
high. The composite is laid up with the principal fibre direction
along the principal-stress trajectory of the running gait.
(ii) Geometry. J-shape with the curved section storing elastic
energy on touchdown and returning it on push-off. Length
$\sim 30\,\text{cm}$ (matched to subject's leg length); curve radius
set by stiffness target (typical stiffness $\sim 50\,\text{N/mm}$ at
the foot).
(iii) Energy return. Target $\sim 70\,\%$ of input energy returned
on push-off (intact human ankle: $\sim 90\,\%$ from Achilles tendon
recoil; current prosthetic blades reach $\sim 70\text{-}80\,\%$).
(iv) Stiffness tuning. Blade thickness varies along the curve to
distribute strain (uniform-strain design). A heavier or faster runner
needs a stiffer blade.
(v) Failure modes. Fatigue delamination at the highest-strain region
of the curve; manufacturing defects (void content $> 1\,\%$ reduces
fatigue life by $\sim 30\,\%$); UV degradation of epoxy matrix over
years. Recommended replacement interval: every $1500\,\text{km}$ of
running, or annually for heavy users.

\begin{exercise}
Design a stent for a coronary artery with a $3\,\text{mm}$ inner
diameter. Specify the material, expanded geometry, and strut
thickness. Defend the choices against the trade-offs (radial force
versus flexibility, drug elution versus bare metal).
\end{exercise}

\paragraph{Discussion.}
Rubric: (i) Material. Cobalt-chrome alloy (L605 or MP35N) or
platinum-chromium for the strut, chosen for high yield strength
(allowing thin struts), radiopacity, and biocompatibility.
Stainless steel acceptable but heavier struts. Bioresorbable
(PLLA, magnesium) considered but recurrent stenosis rates remain
higher than for metal stents.
(ii) Geometry. Open-cell tubular mesh, $3\,\text{mm}$ expanded
diameter, $15\,\text{mm}$ length (typical for a single lesion),
$70$-$90\,\si{\micro\meter}$ strut thickness (thin enough to allow
re-endothelialisation, thick enough to provide radial force
$> 0.1\,\text{N/mm}$ at full expansion).
(iii) Drug coating. Sirolimus or everolimus drug-eluting coating
(polymer matrix releasing drug over $30$-$90$ days) to suppress
neointimal hyperplasia. Drug-eluting stents have $\sim 5\,\%$
restenosis rate at $1\,\text{year}$ vs $\sim 20\,\%$ for bare-metal
stents but require longer dual antiplatelet therapy.
(iv) Trade-off discussion. Thinner struts $\Rightarrow$ faster
endothelialisation, lower thrombosis risk, but lower radial force
(risk of stent recoil). Higher drug dose $\Rightarrow$ more
restenosis suppression but delayed healing and higher late-thrombosis
risk. Current standard balances at $\sim 80\,\si{\micro\meter}$
strut, second-generation drug coating.

\subsection*{Diagnosis and reverse engineering}\pathtag{core}

\begin{exercise}
A patient with knee osteoarthritis reports pain during stair
descent more severe than during stair ascent. Explain the
biomechanical reason and propose one rehabilitation strategy that
addresses it.
\end{exercise}

\paragraph{Solution sketch.}
Stair descent requires the quadriceps to perform eccentric
contraction: the muscle is lengthening while developing tension to
control the body's descent. Eccentric force is $\sim 1.5$ times
maximum isometric force at the same activation, so the
patello-femoral compression rises proportionally. The high-force,
high-flexion, eccentric loading produces the highest patello-femoral
joint compression of common daily activities ($\sim 3500\,\text{N}$
vs $\sim 1500\,\text{N}$ during ascent), the reason pain is worse
in descent. Rehabilitation: quadriceps strengthening through
controlled eccentric exercise (slow squat descents, leg-press
eccentric phase), which both improves the quadriceps' eccentric
capacity and conditions the cartilage's metabolic response. Add
gait modification: descend stairs with the knee less flexed (longer
quadriceps moment arm reduces the required muscle force) and use a
handrail to share load.

\begin{exercise}
A patient's gait analysis shows reduced ankle dorsiflexion moment
during stance. Identify three possible engineering causes (calf
muscle weakness, Achilles tendon shortening, ankle joint
limitation) and propose a diagnostic test for each.
\end{exercise}

\paragraph{Solution sketch.}
The ankle dorsiflexion moment during stance is generated by the
plantar-flexors (gastrocnemius + soleus) acting through the Achilles
tendon. Three possible causes:
(i) Calf muscle weakness. Test: manual muscle test of plantarflexion
strength against resistance; single-leg heel-rise test (normal:
$\sim 25$ repetitions; weak: $< 10$). Cause: post-disuse, motor
neuropathy, prior injury.
(ii) Achilles tendon shortening (contracture). Test: knee-extended
ankle range of motion (normal: $20^{\circ}$ dorsiflexion; contracted:
$<5^{\circ}$); Silfverskiold test (compare knee-flexed and
knee-extended dorsiflexion to localise to gastrocnemius vs soleus).
Cause: post-immobilisation, post-burn, congenital.
(iii) Ankle-joint limitation. Test: passive ankle range with the
knee both flexed and extended; if both are limited and there is bony
end-feel, the limitation is in the joint not the soft tissue; imaging
(weight-bearing CT or fluoroscopy) confirms anterior tibial-talar
impingement or post-traumatic arthrofibrosis.

\begin{exercise}
An athlete's prosthetic running blade fatigue-fractures after
$2$ years of use. Identify three possible engineering causes
(loading exceeds design specification, manufacturing defect,
fatigue-induced delamination of the composite) and propose a
diagnostic measurement.
\end{exercise}

\paragraph{Solution sketch.}
(i) Loading exceeds design specification. The athlete may have
gained mass, increased training mileage (more cycles), or changed
running mechanics (greater impact peaks). Diagnostic: instrument the
remaining blade with strain gauges or compare the failure load to
the design envelope using the manufacturer's specification sheet.
(ii) Manufacturing defect. Voids, fibre misalignment, or
fibre-matrix debonding at the highest-stress region of the curve
can initiate fatigue cracks at well below the nominal cycle
threshold. Diagnostic: C-scan ultrasound of an intact replacement
blade from the same batch; cross-sectional micrograph of the
fractured blade at the crack-initiation site.
(iii) Fatigue-induced delamination. Even a defect-free composite
accumulates microcracks under cyclic loading; the carbon-fibre /
epoxy combination is fatigue-sensitive at $> 60\,\%$ of UTS at
moderate cycle counts. Diagnostic: fractographic SEM of the failed
section to identify fatigue-crack striations vs catastrophic
delamination; count cycles to failure and compare to the published
S-N curve.

\subsection*{Failure analysis}\pathtag{core}

\begin{exercise}
A $75$-year-old patient suffers a hip fracture from a low-energy
fall. Identify the dominant failure mode (osteoporotic fragility
fracture) and the engineering response in modern fracture
treatment.
\end{exercise}

\paragraph{Solution sketch.}
Dominant failure mode: osteoporotic fragility fracture. Mechanism:
trabecular bone density has fallen by $\sim 40\,\%$ from peak
(typical decrement at age 75), so the compressive strength of the
femoral neck has fallen as $\rho^{2.5}$ to $\sim 30\,\%$ of peak,
i.e.\ from $\sim 9\,\text{MPa}$ to $\sim 3\,\text{MPa}$. A low-energy
fall on the hip (subject decelerating $\sim 1\,\text{m}$ fall over
$\sim 0.05\,\text{s}$) produces $\sim 5000\,\text{N}$ across a
$\sim 2\,\text{cm}^2$ femoral neck cross-section, giving
$\sim 25\,\text{MPa}$, well above the degraded strength. Engineering
response: (i) acute, surgical fixation. For displaced intracapsular
fractures: total hip replacement or hemi-arthroplasty (because the
femoral head's blood supply is disrupted). For non-displaced or
extracapsular: dynamic hip screw, intramedullary nail, or
cephalomedullary nail.
(ii) Long-term: bisphosphonate or anti-RANKL anti-resorptive
therapy, calcium and vitamin D supplementation, fall-prevention
intervention (balance training, hip protectors, home modification).
(iii) System-level: outcome metric is the $30$-day mortality
($\sim 8\,\%$ post-hip-fracture, current as of 2024) and the
one-year return-to-baseline-function rate ($\sim 40\,\%$). Modern
care pathways aim for surgery within $48\,\text{h}$ and early
mobilisation to reduce both.

\begin{exercise}
A hip replacement loosens at the bone-implant interface $15$ years
after surgery; revision surgery reveals extensive polyethylene wear
particles in the surrounding tissue. Trace the chain from particle
generation to aseptic loosening.
\end{exercise}

\paragraph{Solution sketch.}
The chain has six steps:
(i) Bearing surface wear. The CoCrMo head sliding on the UHMWPE
acetabular liner generates submicrometre polyethylene particles at
a rate of $\sim 10^{10}$ particles per year for conventional UHMWPE.
Wear rate is $\sim 100\,\si{\micro\meter}/\text{year}$ vertical
penetration.
(ii) Particle migration. Particles diffuse into the periprosthetic
joint space and reach the bone-implant interface through the
effective joint space.
(iii) Macrophage phagocytosis. Periprosthetic macrophages phagocytose
the wear particles. The particle size $0.1$-$1\,\si{\micro\meter}$
is optimal for triggering inflammatory cytokine release (TNF-$\alpha$,
IL-1, IL-6) but particles cannot be enzymatically degraded.
(iv) Osteoclast activation. The cytokine signalling activates
osteoclast precursors; the resulting osteoclast activity resorbs
the periprosthetic bone.
(v) Bone resorption produces gradual radiolucent zones at the
bone-implant interface, visible on X-ray as expanding lucent lines
around the prosthesis.
(vi) Loss of fixation. When the lucent zone exceeds $\sim 2\,\text{mm}$
of continuous resorption around the prosthesis, mechanical
micromotion at the bone-implant interface becomes a loaded gap; the
patient experiences pain on weight-bearing; revision is indicated.
Engineering response: cross-linked UHMWPE (developed in the late
1990s) reduces wear-particle generation by $\sim 90\,\%$ and has
extended the median time to aseptic loosening from $\sim 15$ years
to $\sim 25$ years for primary hip replacements.

\begin{exercise}
A patient develops periprosthetic infection three months after knee
replacement. Identify the most likely engineering source (intra-operative
contamination, haematogenous spread, biofilm formation on the implant
surface) and the engineering responses in revision surgery.
\end{exercise}

\paragraph{Solution sketch.}
At three months post-surgery, the most likely source is
intra-operative contamination of the implant surface during the
primary procedure; the organisms (most commonly coagulase-negative
\emph{Staphylococcus} or \emph{S.\ aureus}) form a biofilm on the
implant surface that resists both host immune clearance and
antibiotic penetration. Haematogenous spread from a distant focus
is the second-most-common source and rises in importance after
the first year post-implantation.

Engineering responses in revision:
(i) Single-stage exchange (current evidence supports this for
low-virulence organisms with good soft-tissue coverage): explant
all hardware, debride the periprosthetic tissue, and re-implant with
antibiotic-impregnated bone cement (typically gentamicin- or
vancomycin-loaded PMMA at $1$-$4\,\text{g}/40\,\text{g}$ cement).
(ii) Two-stage exchange (standard for high-virulence organisms or
poor soft-tissue coverage): explant all hardware; implant an
antibiotic-eluting spacer (PMMA block with antibiotics at high dose);
IV antibiotics for $6$ weeks; confirm infection eradication
(normal CRP, negative cultures); re-implant a permanent prosthesis.
(iii) Adjunctive: optimise nutrition, glycaemic control, smoking
cessation; consider chronic suppressive antibiotic therapy for
patients who cannot tolerate explant.

\subsection*{Judgment}\pathtag{standard}

\begin{exercise}
A press release announces a ``running shoe that reduces knee load
by 30\%''. Identify the engineering measurements that the claim
should be supported by (joint contact force during running, not just
ground reaction force) and write a one-paragraph reply asking the
manufacturer to specify the load measurement.
\end{exercise}

\paragraph{Discussion.}
Rubric (sample reply): Thank you for the announcement. So that
clinicians and biomechanical engineers can interpret the claim,
please specify (i) which knee-load metric was measured: peak
patello-femoral contact force, peak tibio-femoral contact force,
peak knee-extension moment, peak vertical GRF, or impulse over the
stance phase; (ii) the measurement method: in-vivo telemetry from an
instrumented knee prosthesis (the gold standard) or model-based
inverse dynamics; (iii) the sample (sample size, sex, body mass,
running mileage history, running technique), since the reported
benefit may not generalise to runners with different patterns;
(iv) the comparator shoe (the $30\,\%$ reduction is relative to
what); (v) effect size with confidence interval (a $30\,\%$ mean
reduction with a $\pm 40\,\%$ CI is consistent with no effect);
(vi) any pre-registration or independent replication. A reduction
in vertical GRF, without a corresponding reduction in joint contact
force, does not reduce the cartilage compressive stress; muscle
activation typically adapts to keep joint compression near constant
under altered GRF.

\begin{exercise}
A device manufacturer claims that its osteoporosis-prevention
``vibration platform'' produces ``cellular-level bone stimulation''
through low-amplitude high-frequency vibration. Identify the
engineering question the claim must answer (does the imposed strain
reach the bone tissue at sufficient amplitude and frequency to
produce a measurable osteogenic response in vivo) and write a
one-paragraph reply.
\end{exercise}

\paragraph{Discussion.}
Rubric (sample reply): The mechanostat literature
\cite[p.~74]{paper:v6c06-frost-1987} suggests that bone formation is
triggered above a strain threshold ($\sim 1000$ to $2000\,$microstrain).
Please specify (i) the in-vivo tissue strain at the bone of interest
(femoral neck, lumbar vertebra) produced by your device at the
manufacturer-specified vibration parameters: this requires modelling
or direct strain-gauge measurement, not just the platform's surface
acceleration; (ii) the strain rate (frequency content): the response
curve is dominated by frequencies in the $10$-$50\,\text{Hz}$ range,
with little effect above $100\,\text{Hz}$; (iii) the
placebo-controlled randomised-clinical-trial evidence of bone
mineral density gain at one year, with the magnitude expressed as a
$z$-score or a $T$-score change; (iv) independent replication in a
population matched to your indicated use (post-menopausal women,
elderly men). Without the in-vivo strain estimate, the ``cellular
stimulation'' claim cannot be evaluated.

\begin{exercise}
A reader argues that biomechanics is ``not really engineering''
because the body is too complex to model accurately. Identify two
engineering achievements of biomechanics (joint replacement design,
gait analysis for orthotic prescription) and one engineering
limitation (patient-specific calibration is rarely done) and write
a one-paragraph reply.
\end{exercise}

\paragraph{Discussion.}
Rubric (sample reply): The objection conflates ``cannot be modelled
perfectly'' with ``cannot be engineered''. The Charnley low-friction
arthroplasty was a designed device with quantified material
properties, fatigue testing, and outcome metrics, and has been
implanted in tens of millions of patients with twenty-year revision
rates near $10\,\%$ \cite[ch.~14]{text:mow-huiskes-orthopaedic}.
Modern gait analysis takes a few-hundred-point force-plate trace,
runs inverse dynamics, and produces a joint-moment-time curve that
informs orthotic and surgical decisions for cerebral palsy,
post-stroke, and amputation patients. The valid limitation is
specificity: most biomechanical models use population-average
segment masses, inertias, and muscle parameters, while the
between-subject variation is $\pm 20\,\%$ or more. Patient-specific
imaging-based modelling is the active research frontier, not the
clinical standard. The right summary is that biomechanics is
engineering at the level of population-average design, and is
moving toward patient-specific design as imaging and computation
costs fall.

\begin{exercise}
A textbook describes the heart as ``a pump''. Identify two ways in
which the pump metaphor is accurate (pressure generation, flow
direction) and two ways in which it is misleading (the heart's
output is pulsatile, and the pump itself is built from a single cell
type that is also electrically active).
\end{exercise}

\paragraph{Discussion.}
Rubric: Accurate aspects: (i) The heart generates pressure
($\sim 100\,\text{mmHg}$ mean arterial) from mechanical work
($\sim 1\,\text{W}$ at rest) on a working fluid (blood); the
PV-loop area is the stroke work, the same diagram used for a
positive-displacement compressor. (ii) The heart enforces flow
direction through one-way valves (mitral, tricuspid, aortic,
pulmonary), the same engineering device as the check valves in a
mechanical pump.
Misleading aspects: (i) The output is pulsatile, not continuous,
which has consequences for arterial wave propagation (Moens-Korteweg)
and for the design of mechanical heart assist devices: most
continuous-flow LVADs produce non-physiological pulse pressure with
downstream consequences (gut bleeding, aortic insufficiency).
(ii) The cardiac pump is built from cells (cardiomyocytes) that are
both the actuators and the signal-processing elements: the
electrical conduction system is built from the same cell type as
the contractile myocardium, only with different gene expression.
Pumps engineered by humans separate the actuator from the
controller; biology rarely does.

\begin{exercise}
A clinical-trial protocol proposes to use pulse-wave velocity as the
primary endpoint for a study of a new anti-hypertensive drug.
Identify two engineering advantages of the choice (mechanistic
linkage to arterial stiffness, non-invasive measurement) and two
limitations (high inter-observer variability, dependence on
ambient blood pressure at the time of measurement). Write a
one-paragraph reply asking the protocol's authors to specify how
they will control the limitations.
\end{exercise}

\paragraph{Discussion.}
Rubric (sample reply): PWV is mechanistically linked to arterial
wall stiffness through the Moens-Korteweg equation, and is
non-invasively measurable in a $\sim 10$-minute outpatient visit,
both genuine engineering advantages. Limitations: (i) PWV scales
with the square root of the wall modulus but also depends on the
mean blood pressure at the time of measurement; a fall in PWV could
reflect drug-induced arterial relaxation or simply a transient
lower BP at the visit. Please specify how you will control for
ambient BP (concurrent BP measurement, fixed time-of-day,
medication-free interval). (ii) Inter-observer variability for
carotid-femoral PWV is $\sim 0.5\,\text{m/s}$ ($\sim 7\,\%$
relative) for trained operators and worse for untrained;
inter-device variability can reach $1\,\text{m/s}$. Please specify
operator training, device standardisation, and central-reading
protocol. With these controls, PWV is a defensible primary endpoint
for an arterial-stiffness mechanism trial; without them, the
endpoint is too noisy to deliver a useful effect estimate.
